% ============================================================================
%  DDS-AMQP 1.0 — ZeroDDS Vendor Specification.
%
%  Build:  tectonic main.tex
%  Output: main.pdf  (renamed to dds-amqp-1.0-beta1.pdf in CI release jobs)
% ============================================================================

\documentclass{../styles/zerodds-vendor-spec}

\title{DDS-AMQP\\\large Bridge Specification for OMG DDS over OASIS AMQP 1.0}
\specshorttitle{DDS-AMQP 1.0}
\specversion{zerodds/v1.0.0-beta1}
\specstatus{Beta}
\specdate{2026-04-29}
\speceditors{Sandra Ke{\ss}ler}

\begin{document}

\maketitle

\tableofcontents
\clearpage

% ============================================================================
%  Status of this Document.
% ============================================================================
\chapter*{Status of this Document}
\addcontentsline{toc}{chapter}{Status of this Document}

This is a Beta release of the DDS-AMQP Vendor Specification, version
\texttt{zerodds/v1.0.0-beta1}, published by the ZeroDDS Project. It is
made available for public review and feedback prior to Final
publication.

\begin{specwarning}
\textbf{Wire-format change vs. earlier beta1 drafts.} This
revision changes the on-the-wire contract in several ways
relative to earlier \texttt{beta1} working drafts:
\texttt{message-id} is now a 24-byte per-sample identifier
(was: 16-byte \texttt{InstanceHandle\_t});
\texttt{group-id} is now a SHA-256 hex digest of the XCDR2
KeyHash (was: SOLIDUS-separated lexical key string);
the SASL PLAIN, ANONYMOUS, and SCRAM-SHA-256 IdentityTokens
now use \texttt{zerodds:Auth:*} class\_ids (was:
\texttt{DDS:Auth:PKI-DH:1.0} for all);
the \texttt{TopicMapping.dds\_partition} field is now a
\texttt{sequence<string>} (was: scalar string);
the \texttt{bridge\_id}/\texttt{bridge\_hop\_cap} fields have
moved from \texttt{TopicMapping} to \texttt{AmqpEndpointConfig}
and \texttt{AmqpBridgeConfig}.
The \texttt{beta} status exists precisely for changes of this
shape; first ratified release will be \texttt{v1.0.0}.
\end{specwarning}

\paragraph{Versioning Scheme.}
This specification follows the SemVer-2.0 versioning convention:
\texttt{zerodds/vMAJOR.MINOR.PATCH} with optional pre-release
suffix \texttt{-betaN} or \texttt{-rcN}. Examples:

\begin{itemize}
  \item \texttt{zerodds/v1.0.0-beta1} — current Beta iteration.
  \item \texttt{zerodds/v1.0.0-rc1} — first Release Candidate.
  \item \texttt{zerodds/v1.0.0} — Final release.
  \item \texttt{zerodds/v1.0.1} — first editorial errata.
\end{itemize}

Git-tag names mirror the version with the prefix
\texttt{spec-dds-amqp-} and replace the slash with a hyphen
(e.g.\ \texttt{spec-dds-amqp-v1.0.0-beta1},
\texttt{spec-dds-amqp-v1.0.0}, \texttt{spec-dds-amqp-v1.0.1}).

\paragraph{Date Field.}
The date printed on the title page is the date of the last
\emph{editorial} change to this specification text. It is
\rfckw{NOT} the build-date of the PDF artefact; build-date is
recorded in the PDF metadata only.

This specification defines a normative bridge between the OMG Data
Distribution Service for Real-Time Systems (DDS) and the OASIS AMQP
1.0 messaging protocol. It does not modify, profile, or constrain the
underlying DDS or AMQP standards; it specifies the mapping between
them and the conformance requirements an implementation \rfckw{MUST}
satisfy.

\paragraph{Subset Conformance.}
This specification does not require an implementation to claim
\emph{full} OMG DDS or OASIS AMQP conformance certification. An
implementation \rfckw{MAY} build a DDS-AMQP-conformant bridge using
a partial DDS or AMQP stack, provided that every DDS feature and
every AMQP feature \emph{actually used} on either side of the
bridge is itself implemented in conformance with the relevant cited
normative source.

In practice this means:

\begin{itemize}
  \item The AMQP-side of a DDS-AMQP implementation \rfckw{MUST}
        speak AMQP-1.0 wire-protocol that is byte-conformant with
        \omgref{amqp-1.0-transport}{} and
        \omgref{amqp-1.0-types}{} for the operations it
        performs (no proprietary AMQP framing).
  \item The DDS-side \rfckw{MUST} produce DCPS-conformant
        DataReader/DataWriter behaviour
        (\omgref{OMG DDS}{\S\,2.2.2}) and XCDR2-conformant
        wire-encoding (\omgref{OMG X-Types}{\S\,7.4.3}) for the
        topic types it serves.
  \item Specifications layered on top of DCPS (e.g.\ DDS-RPC
        request-reply patterns per \omgref{OMG DDS-RPC}{1.0}) are
        transparently transportable in Pass-Through mode
        (\S\,\ref{sec:psm-passthrough}) without being separately
        certified by the DDS-AMQP implementation. An RPC-aware
        mapping that lifts the DDS-RPC correlation header into
        AMQP message properties is given in Annex~D (normative
        when \texttt{rpc\_aware = true}).
  \item The full \omgref{AMQP-Management}{1.0} protocol is
        out-of-scope for this specification (see
        \S\,\ref{sec:out-of-scope}). For operational visibility,
        this specification defines its \emph{own} minimal,
        read-only management surface on the Endpoint Profile —
        \texttt{\$catalog} and \texttt{\$metrics}
        (\S\,\ref{sec:management-interface}) and the security
        \texttt{\$audit} channel (\S\,\ref{sec:audit-channel}) —
        which is independent of AMQP-Management-1.0 and does not
        constitute a partial implementation of it.
\end{itemize}

This subset-conformance posture is consistent with the bridge
specification's role as a mapping document; it neither imposes
nor relaxes the OMG/OASIS conformance bars on either side, but
states which subsets are touched.

\paragraph{Interoperability Constraint.}
Subset Conformance does \emph{not} permit an implementation to
silently skip features within a profile it has declared. The
profiles defined in Chapter~\ref{ch:conformance} (Endpoint,
Bridge, Codec, Codec-Lite) are \emph{sealed conformance units}:
a claim of conformance to a profile \rfckw{MUST} satisfy every
item of that profile's conformance clause without exception.

What Subset Conformance \emph{does} permit is the choice of
\emph{which profile} to claim — for example, a deployment that
uses only the Pass-Through encoding mode \rfckw{MAY} declare
Codec-Lite Profile (\S\,\ref{sec:profile-codec-lite}) instead of
the full Codec Profile. Two implementations claiming the same
profile \rfckw{SHALL} be byte-interoperable in the operations
that profile defines; mixing implementations that claim
different profiles \rfckw{SHALL} require an interoperability
analysis based on the profile-intersection.

The terms \rfckw{MUST}, \rfckw{MUST NOT}, \rfckw{SHALL},
\rfckw{SHALL NOT}, \rfckw{SHOULD}, \rfckw{SHOULD NOT}, \rfckw{MAY},
\rfckw{MAY NOT}, \rfckw{REQUIRED}, \rfckw{RECOMMENDED}, and
\rfckw{OPTIONAL} are to be interpreted as described in IETF RFC 2119
and RFC 8174.

% ============================================================================
%  Chapter 1 — Scope.
% ============================================================================
\chapter{Scope}

\section{Purpose}

This specification defines a normative mapping between the OMG Data
Distribution Service (DDS) data model and the OASIS Advanced Message
Queuing Protocol (AMQP) version 1.0. The mapping enables AMQP
producers and consumers to interoperate with DDS publishers and
subscribers without an intermediate broker, by means of a
DDS-AMQP-conformant endpoint.

Four conformance profiles are defined (see \S\,\ref{ch:conformance}):

\begin{itemize}
  \item the \emph{Endpoint Profile}, applicable to implementations
        that act as a server-side AMQP endpoint and expose DDS topics
        to AMQP clients;
  \item the \emph{Bridge Profile}, applicable to implementations that
        act as an AMQP client and bridge between an external AMQP
        broker (e.g.\ RabbitMQ, Apache Qpid Broker-J, Solace
        PubSub+, Microsoft Azure Service Bus) and DDS;
  \item the \emph{Codec Profile}, applicable to implementations of
        the AMQP type-system, performatives, and message-section
        codecs only, without connection or session lifecycle support;
  \item the \emph{Codec-Lite Profile} (\S\,\ref{sec:profile-codec-lite}),
        a strict subset of Codec for deployments that need only the
        primitive type codecs and the \texttt{data} body section.
\end{itemize}

\section{Relationship to Other DDS Bridge Specifications}\label{sec:relationship-other-bridges}

DDS-AMQP is one of several DDS bridge or transport specifications.
This section is informative and clarifies when a deployment
should select DDS-AMQP versus another bridge.

\begin{longtable}{p{3.4cm}p{3.4cm}p{5.2cm}}
  \toprule
  \textbf{Specification} & \textbf{Primary use-case} & \textbf{Relation to DDS-AMQP} \\
  \midrule
  \endhead
  \omgref{OMG DDS-Web}{1.0}    & Web-browser-to-DDS via REST/HTTP, RESTful resource model  & complementary; DDS-Web targets stateless HTTP request/response semantics. It is \emph{not} a functional replacement for streaming pub/sub. Browser-side streaming pub/sub is out of scope for this revision. \\
  \omgref{OMG DDS-RPC}{1.0}    & request/response patterns layered on DCPS                  & supported in two modes (\S\,\ref{sec:rpc-interop-caveat} and Annex~\ref{annex:rpc-mapping}). \emph{Pass-Through mode} (\texttt{rpc\_aware = false}, default) transports DDS-RPC samples transparently — the DDS-RPC correlation header travels inside the XCDR2 body verbatim, and DDS-aware consumers are required at both ends. \emph{RPC-aware mode} (\texttt{rpc\_aware = true}) lifts the DDS-RPC correlation into AMQP \texttt{message-id}, \texttt{correlation-id}, and \texttt{reply-to} per Annex~D, with the reply-validation semantics defined in \S\,D.4. \\
  \omgref{OMG DDS-XRCE}{1.0}   & extremely-resource-constrained-environments (sensors, MCUs) & alternative for tiny embedded targets that cannot host an AMQP stack; DDS-AMQP targets server-class and edge-gateway-class nodes. \\
  MQTT v3.1.1 / v5.0           & lightweight pub/sub for IoT / mobile                       & disjoint wire-protocol family. Deployments that need MQTT connectivity \rfckw{SHOULD} use a dedicated MQTT-DDS bridge (out of scope for this specification); DDS-AMQP does not bridge MQTT directly. MQTT is the de-facto standard for resource-constrained IoT devices; DDS-AMQP is positioned for AMQP-1.0-native enterprise messaging fabrics (RabbitMQ, Solace, Azure Service Bus). An Endpoint or Bridge Profile implementation \rfckw{MAY} coexist with an MQTT-DDS bridge on the same DDS domain. Forward-loop prevention is normatively specified in \S\,\ref{sec:bridge-coexistence} (\texttt{dds:bridge-id} stamping, drop-on-self-tag, hop-count cap) and applies to every Endpoint Profile and Bridge Profile implementation that forwards samples between DDS and AMQP. \\
  AMQP-WebSockets-1.0          & AMQP-1.0 over WebSocket transport                         & out of scope for this revision (\S\,\ref{sec:out-of-scope}). Not part of v1.0.0-beta1 conformance; addressed in a future revision. \\
  AMQP 0-9-1 / 0-10            & legacy AMQP wire-protocols                                 & disjoint protocol family. The relevant distinction is the wire protocol on the hop adjacent to the DDS-AMQP endpoint: a broker speaking native AMQP-1.0 on that hop (RabbitMQ 3.8+ with native AMQP-1.0, Apache Qpid Broker-J, Apache ActiveMQ Artemis, Solace PubSub+, Microsoft Azure Service Bus) is a conformant counterparty, regardless of what other protocols it speaks elsewhere. Topologies that interpose a 0-9-1$\leftrightarrow$1.0 protocol-translator hop between the endpoint and the AMQP-1.0 broker are out of scope (\S\,\ref{sec:out-of-scope}); this specification does not characterise their behaviour. \\
  DDSI-RTPS over TCP / WAN     & native DDS wire protocol routed across network boundaries  & alternative to a protocol-bridge entirely. Deployments that control both endpoints (DDS native on each side) \rfckw{SHOULD} consider native DDSI-RTPS routing (TCP transport, WAN-aware discovery services) before introducing a protocol-bridge. DDS-AMQP is the right choice when one side of the boundary speaks AMQP-1.0 (broker, cloud-messaging) and a wire-format break is unavoidable. \\
  \bottomrule
\end{longtable}

A typical deployment pattern uses DDS-AMQP for edge-gateway-to-cloud
federation (durable, broker-mediated or direct-embed) and DDS-Web
for browser-side request/response interactions; for browser-side
streaming pub/sub the WebSockets-binding caveat above applies.

\subsection*{RPC Interoperability Caveat}\label{sec:rpc-interop-caveat}

A DDS-RPC sample transported in Pass-Through mode is opaque to
non-DDS AMQP consumers: the AMQP \texttt{reply-to} and
\texttt{correlation-id} properties are \emph{not} populated by
the DDS-RPC correlation header (which lives inside the XCDR2 body).
Native AMQP-RPC clients (e.g.\ those using the AMQP-1.0 dynamic-
reply-to convention from
\omgref{amqp-1.0-messaging}{\S\,3.2.4}) will therefore not
interoperate with DDS-RPC traffic in Pass-Through mode.

For interoperability with native AMQP-RPC, the topic
configuration \rfckw{MUST} set \texttt{rpc\_aware = true}
(Annex~A); the endpoint then lifts the DDS-RPC correlation
header into AMQP properties as defined in Annex~D, making the
exchange a first-class AMQP-RPC pattern.

\section{Out of Scope}\label{sec:out-of-scope}

The following are explicitly outside the scope of this specification:

\begin{itemize}
  \item AMQP broker functionality (queue persistence, exchange
        routing, dead-letter handling, message TTL, and similar
        features traditionally provided by AMQP brokers).
        Implementations that require these features \rfckw{SHOULD}
        deploy a conformant broker between the AMQP client and the
        DDS-AMQP endpoint, in the configuration described
        in~\S\,\ref{sec:topology-sidecar}.
  \item The full AMQP Management protocol (AMQP-Management-1.0).
        Note that this specification provides its own minimal,
        read-only management surface — \texttt{\$catalog} and
        \texttt{\$metrics} (\S\,\ref{sec:management-interface})
        and the \texttt{\$audit} channel
        (\S\,\ref{sec:audit-channel}). That surface is defined
        natively in this document and is not a profile or subset
        of AMQP-Management-1.0.
  \item AMQP Web-Sockets binding (AMQP-WebSockets-1.0); this binding
        \rfckw{MAY} be added in a future revision.
  \item Modifications to the OMG DDS data model or behaviour. Any
        DDS data type, QoS policy, or behaviour referenced by this
        specification is governed by its OMG-normative source.
  \item AMQP 0-9-1 or AMQP 0-10 protocol versions; this specification
        is bound exclusively to OASIS AMQP 1.0.
\end{itemize}

% ============================================================================
%  Chapter 2 — Conformance.
% ============================================================================
\chapter{Conformance}\label{ch:conformance}

This specification defines four conformance profiles. An
implementation \rfckw{MAY} claim conformance to one or more profiles.
A conformance claim \rfckw{MUST} cite this specification by full
identifier (\texttt{DDS-AMQP, zerodds/v1.0.0-beta1} or later) and
enumerate every profile to which conformance is claimed.

\section{Endpoint Profile}\label{sec:profile-endpoint}

\begin{conformance}{Endpoint Profile}
An implementation conforms to the Endpoint Profile if and only if all
of the following are true:
\begin{enumerate}
  \item It supports the AMQP type-system as specified
        in~\S\,\ref{sec:pim-type-mapping}.
  \item It accepts incoming AMQP connections in accordance
        with~\S\,\ref{sec:topology-direct} and the AMQP 1.0
        connection-state model defined in
        \omgref{amqp-1.0-transport}{\S\,2.4}.
  \item It accepts both Sender and Receiver link attachments
        as specified in~\S\,\ref{sec:pim-settlement}, and translates
        them bidirectionally into DDS DataWriter and DataReader
        operations.
  \item It implements address resolution as specified
        in~\S\,\ref{sec:pim-address-resolution}.
  \item It supports at least the Pass-Through encoding mode
        defined in~\S\,\ref{sec:psm-passthrough}, and \rfckw{SHOULD}
        support the JSON encoding mode of~\S\,\ref{sec:psm-json}.
  \item It supports the SASL PLAIN, ANONYMOUS, and EXTERNAL
        mechanisms as specified in~\S\,\ref{sec:security-sasl}.
  \item It \rfckw{SHALL NOT} offer SASL PLAIN to a client whose
        underlying transport is unencrypted
        (\S\,\ref{sec:security-sasl-plain-tls}). A client SASL-init
        frame carrying \texttt{PLAIN} on an unencrypted connection
        \rfckw{SHALL} be rejected with SASL outcome \texttt{auth},
        and the connection \rfckw{SHALL} be closed.
  \item It enforces the No-Bypass guarantee for DDS-Security
        identities (\S\,\ref{sec:security-no-bypass}) and
        per-link Governance resolution
        (\S\,\ref{sec:per-link-governance}).
  \item It exposes the operational management surface defined
        in~\S\,\ref{sec:management-interface} (the
        \texttt{\$catalog} and \texttt{\$metrics} addresses) and
        the security audit channel defined
        in~\S\,\ref{sec:audit-channel} (the \texttt{\$audit}
        address).
  \item It implements the bridge-coexistence rules
        (\S\,\ref{sec:bridge-coexistence}) for every sample it
        forwards between DDS and AMQP.
  \item It satisfies the conformance test cases enumerated in
        Annex~C, \S\,C.1.
\end{enumerate}
\end{conformance}

\section{Bridge Profile}\label{sec:profile-bridge}

\begin{conformance}{Bridge Profile}
An implementation conforms to the Bridge Profile if and only if all
of the following are true:
\begin{enumerate}
  \item It supports the AMQP type-system as specified
        in~\S\,\ref{sec:pim-type-mapping}.
  \item It establishes outgoing AMQP connections to an external
        broker as specified in~\S\,\ref{sec:topology-sidecar}, and
        attaches Sender and Receiver links translating
        bidirectionally with DDS DataWriters and DataReaders.
  \item It implements address resolution as specified
        in~\S\,\ref{sec:pim-address-resolution}.
  \item It supports the SASL PLAIN, ANONYMOUS, and EXTERNAL
        mechanisms as specified in~\S\,\ref{sec:security-sasl}
        for the outbound connection to the broker.
  \item It \rfckw{SHALL NOT} initiate a SASL PLAIN authentication
        over an unencrypted transport
        (\S\,\ref{sec:security-sasl-plain-tls}).
  \item It honours the dual-identity model of
        \S\,\ref{sec:bridge-dual-identity}: broker-side
        SASL credentials \rfckw{SHALL NOT} be conflated with the
        bridge's DDS-Security IdentityToken.
  \item It performs reconnect with exponential back-off after
        connection loss as specified in~\S\,\ref{sec:reconnect}.
  \item It supports at least the Pass-Through encoding mode of
        \S\,\ref{sec:psm-passthrough}, and \rfckw{SHOULD} support
        the JSON encoding mode of~\S\,\ref{sec:psm-json} for
        interoperability with non-DDS-aware AMQP clients
        (matching the Endpoint Profile recommendation).
  \item It exposes the operational management surface defined
        in~\S\,\ref{sec:management-interface} (the
        \texttt{\$catalog} and \texttt{\$metrics} addresses) and
        the security audit channel defined
        in~\S\,\ref{sec:audit-channel} (the \texttt{\$audit}
        address) for the DDS-side surface of the bridge; the
        broker-side surface is governed by the colocated
        broker.
  \item It implements the bridge-coexistence rules
        (\S\,\ref{sec:bridge-coexistence}) for every sample it
        forwards between DDS and AMQP.
  \item It satisfies the conformance test cases enumerated in
        Annex~C, \S\,C.2.
\end{enumerate}
\end{conformance}

\section{Codec Profile}\label{sec:profile-codec}

\begin{conformance}{Codec Profile}
An implementation conforms to the Codec Profile if and only if all
of the following are true:
\begin{enumerate}
  \item It implements all primitive type codecs of OASIS AMQP 1.0
        \omgref{amqp-1.0-types}{\S\,1.6.1--\S\,1.6.21}.
  \item It implements the compound-type codecs (list, map, array)
        of \omgref{amqp-1.0-types}{\S\,1.6.22--\S\,1.6.24}.
  \item It implements all nine performatives of
        \omgref{amqp-1.0-transport}{\S\,2.7}: \texttt{open},
        \texttt{begin}, \texttt{attach}, \texttt{flow},
        \texttt{transfer}, \texttt{disposition}, \texttt{detach},
        \texttt{end}, \texttt{close}.
  \item It implements all message sections of
        \omgref{amqp-1.0-messaging}{\S\,3}: \texttt{header},
        \texttt{delivery-annotations}, \texttt{message-annotations},
        \texttt{properties}, \texttt{application-properties}, the
        body section variants \texttt{data} / \texttt{amqp-sequence}
        / \texttt{amqp-value}, and \texttt{footer}.
  \item It satisfies the conformance test cases enumerated in
        Annex~C, \S\,C.3.
\end{enumerate}
\end{conformance}

\section{Codec-Lite Profile}\label{sec:profile-codec-lite}

\begin{conformance}{Codec-Lite Profile}
An implementation conforms to the Codec-Lite Profile if and only
if all of the following are true:
\begin{enumerate}
  \item It implements all primitive type codecs of OASIS AMQP 1.0
        \omgref{amqp-1.0-types}{\S\,1.6.1--\S\,1.6.21}.
  \item It implements the AMQP \texttt{data} body-section codec
        (\omgref{amqp-1.0-messaging}{\S\,3.2.6}).
  \item It implements all nine performatives of
        \omgref{amqp-1.0-transport}{\S\,2.7}.
  \item It satisfies the conformance test cases in Annex~C,
        \S\,C.4 (Codec-Lite-Profile Tests).
\end{enumerate}
\end{conformance}

The Codec-Lite Profile is a strict subset of the Codec Profile.
It is intended for Endpoint or Bridge deployments that operate
exclusively in Pass-Through and/or JSON encoding modes, where
the AMQP body is always a single \texttt{data} section and no
list/map/array codecs are exercised.

\section{Combination of Profiles}

An implementation \rfckw{MAY} claim conformance to multiple
profiles simultaneously.

\subsection{Codec-Profile Implication for Endpoint and Bridge}\label{sec:codec-implication}

An Endpoint Profile or Bridge Profile implementation
\rfckw{SHALL} additionally claim conformance to either the
Codec Profile (\S\,\ref{sec:profile-codec}) or the Codec-Lite
Profile (\S\,\ref{sec:profile-codec-lite}), as follows:

\begin{longtable}{p{4.0cm}p{8.0cm}}
  \toprule
  \textbf{Encoding modes supported} & \textbf{Required codec profile} \\
  \midrule
  \endhead
  Pass-Through only         & Codec-Lite \\
  Pass-Through + JSON       & Codec-Lite \\
  Pass-Through + JSON + AMQP-Native & Codec (full) \\
  \bottomrule
\end{longtable}

Each profile is a sealed conformance unit with fixed guarantees;
unlike earlier drafts, an Endpoint or Bridge does not silently
skip codec items. The implementation declares which codec
profile it implements, and that profile's full requirements
apply without exception.

\subsection{Hybrid Profile Combination}

A combined Endpoint+Bridge implementation realises the Hybrid
topology of~\S\,\ref{sec:topology-hybrid}. A claim of conformance
to both profiles \rfckw{MUST} satisfy the union of all test cases of
Annex~C, \S\,C.1 and~\S\,C.2.

% ============================================================================
%  Chapter 3 — Normative References.
% ============================================================================
\chapter{Normative References}

The following documents, in whole or in part, are normatively
referenced in this specification and are indispensable for its
application. For dated references, only the edition cited applies.
For undated references, the latest edition of the referenced
document (including any amendments) applies.

\begin{longtable}{p{3.2cm}p{2.0cm}p{8.5cm}}
  \toprule
  \textbf{Identifier} & \textbf{Version} & \textbf{Title} \\
  \midrule
  \endhead
  OMG DDS              & 1.4         & Data Distribution Service for Real-Time Systems\\
  OMG DDSI-RTPS        & 2.5         & Real-Time Publish-Subscribe Wire Protocol\\
  OMG X-Types          & 1.3         & Extensible and Dynamic Topic Types for DDS\\
  OMG DDS-Security     & 1.2         & DDS Security Specification\\
  OMG DDS-XML          & 1.0         & XML Configuration of DDS Resources\\
  OMG IDL              & 4.2         & Interface Definition Language\\
  OASIS AMQP overview  & 1.0         & Advanced Message Queuing Protocol --- Overview\\
  OASIS AMQP types     & 1.0         & Advanced Message Queuing Protocol --- Type System\\
  OASIS AMQP transport & 1.0         & Advanced Message Queuing Protocol --- Transport\\
  OASIS AMQP messaging & 1.0         & Advanced Message Queuing Protocol --- Messaging\\
  OASIS AMQP security  & 1.0         & Advanced Message Queuing Protocol --- Security\\
  IETF RFC 2119        & ---         & Key words for use in RFCs to Indicate Requirement Levels\\
  IETF RFC 8174        & ---         & Ambiguity of Uppercase vs Lowercase in RFC 2119 Key Words\\
  IETF RFC 4422        & ---         & Simple Authentication and Security Layer (SASL)\\
  IETF RFC 4616        & ---         & The PLAIN Simple Authentication and Security Layer Mechanism\\
  IETF RFC 4648        & ---         & The Base16, Base32, and Base64 Data Encodings\\
  IETF RFC 5802        & ---         & Salted Challenge Response Authentication Mechanism (SCRAM)\\
  IETF RFC 8259        & ---         & The JavaScript Object Notation (JSON) Data Interchange Format\\
  IETF RFC 8446        & ---         & The Transport Layer Security (TLS) Protocol Version 1.3\\
  IEEE 754             & 2008        & Standard for Floating-Point Arithmetic\\
  ISO/IEC 9899         & 2018        & Programming languages --- C\\
  \bottomrule
\end{longtable}

\section*{Informative References}

The following documents are referenced informatively only (e.g.\
in the comparison table of \S\,\ref{sec:relationship-other-bridges})
and are not normative for DDS-AMQP conformance.

\begin{longtable}{p{3.4cm}p{2.0cm}p{8.5cm}}
  \toprule
  \textbf{Identifier} & \textbf{Version} & \textbf{Title} \\
  \midrule
  \endhead
  OMG DDS-Web          & 1.0         & Web-Enabled DDS \\
  OMG DDS-RPC          & 1.0         & DDS Remote Procedure Call \\
  OMG DDS-XRCE         & 1.0         & DDS for Extremely Resource Constrained Environments \\
  AMQP-Management      & 1.0         & OASIS AMQP Management protocol (referenced for cross-comparison; out-of-scope for this specification per \S\,\ref{sec:management-interface}) \\
  \bottomrule
\end{longtable}

% ============================================================================
%  Chapter 4 — Terms and Definitions.
% ============================================================================
\chapter{Terms and Definitions}

For the purposes of this document, the following terms and
definitions apply.

\begin{description}[leftmargin=2cm,style=nextline]
  \item[Endpoint] An implementation conforming to the Endpoint
        Profile (\S\,\ref{sec:profile-endpoint}). The endpoint is
        addressed by AMQP clients via a TCP listening socket and
        translates incoming AMQP traffic into DDS operations and
        outgoing DDS traffic into AMQP transfers.
  \item[Bridge] An implementation conforming to the Bridge Profile
        (\S\,\ref{sec:profile-bridge}). The bridge connects to an
        external AMQP broker as a client and translates between the
        broker and DDS.
  \item[Codec] An implementation conforming to the Codec Profile
        (\S\,\ref{sec:profile-codec}). The codec implements the
        AMQP type system, performatives, and message sections
        without connection-lifecycle support.
  \item[Address] A string identifier used by AMQP to designate a
        node within a container; see
        \specref{this specification}{\S\,\ref{sec:pim-address-resolution}}.
  \item[Topic] A DDS-defined named resource carrying samples of a
        single registered type; see \omgref{OMG DDS}{\S\,2.2.2.1}.
  \item[Settlement Mode] An AMQP delivery-state property
        determining whether the receiver is required to acknowledge
        receipt of a transfer; see
        \omgref{amqp-1.0-transport}{\S\,2.6.4}.
  \item[Container] An AMQP-1.0 abstraction representing the
        application instance attaching to a connection; see
        \omgref{amqp-1.0-transport}{\S\,2.4.1}.
  \item[Link] A unidirectional flow of messages between two
        terminuses identified by a name; see
        \omgref{amqp-1.0-transport}{\S\,2.6.1}.
  \item[Body Encoding Mode] One of the three bidirectional
        message-body translation modes defined
        in~\S\,\ref{ch:psm}: Pass-Through, JSON, or AMQP-Native.
  \item[XCDR2] The Extended Common Data Representation, version 2,
        wire format for DDS samples; see
        \omgref{OMG X-Types}{\S\,7.4.3}.
  \item[Address Catalog] An endpoint-exported catalog of available
        topic mappings, as defined
        in~\S\,\ref{sec:pim-discovery-bridging}.
\end{description}

% ============================================================================
%  Chapter 5 — Symbols and Abbreviations.
% ============================================================================
\chapter{Symbols and Abbreviations}

\begin{description}[leftmargin=2cm]
  \item[AMQP] Advanced Message Queuing Protocol
  \item[BNF] Backus-Naur Form
  \item[DDS] Data Distribution Service
  \item[DDSI-RTPS] DDS Interoperability Real-Time Publish-Subscribe
  \item[DCPS] Data-Centric Publish-Subscribe
  \item[IDL] Interface Definition Language
  \item[QoS] Quality of Service
  \item[PIM] Platform-Independent Model
  \item[PSM] Platform-Specific Model
  \item[SASL] Simple Authentication and Security Layer
  \item[SCRAM] Salted Challenge Response Authentication Mechanism
  \item[SPDP] Simple Participant Discovery Protocol
  \item[SEDP] Simple Endpoint Discovery Protocol
  \item[TLS] Transport Layer Security
  \item[XCDR2] Extended Common Data Representation, version 2
\end{description}

% ============================================================================
%  Chapter 6 — Overview.
% ============================================================================
\chapter{Overview (Informative)}

\section{Reference Architecture}

The DDS-AMQP bridge is defined in three reference topologies. An
implementation conforming to the Endpoint Profile realises the
\emph{Direct-Embed} topology of~\S\,\ref{sec:topology-direct}; an
implementation conforming to the Bridge Profile realises the
\emph{Sidecar} topology of~\S\,\ref{sec:topology-sidecar}; the
\emph{Hybrid} topology of~\S\,\ref{sec:topology-hybrid} combines
both within a single deployment.

\subsection{Direct-Embed Topology}\label{sec:topology-direct}

In the Direct-Embed topology, the DDS-AMQP endpoint accepts AMQP
client connections directly. There is no intermediate broker; the
endpoint terminates SASL and TLS, manages the AMQP connection,
session, and link state, and translates link transfers
bidirectionally with DDS data flow. A single endpoint may serve
many concurrent AMQP clients and may participate in any number of
DDS domains simultaneously; the AMQP-side multi-domain addressing
is governed by the \texttt{domain://}-URL-form of the address
resolution rules in~\S\,\ref{sec:pim-address-resolution}.

Bidirectional translation between AMQP transfers and DDS samples
\rfckw{SHALL} honour the AMQP settlement-mode mapping defined
in~\S\,\ref{sec:pim-settlement} (\texttt{snd-settle-mode} and
\texttt{rcv-settle-mode} on the \texttt{attach} performative
correspond to DDS \texttt{RELIABILITY} per the bidirectional
table in that section).

\begin{center}
\resizebox{\textwidth}{!}{%
\begin{tikzpicture}[
  node distance=0.9cm and 0.9cm,
  every node/.style={font=\small},
  box/.style={rectangle,draw=zerodds-primary,rounded corners=2pt,
              minimum width=2.0cm,minimum height=0.9cm,align=center,
              fill=zerodds-light},
  ddsbox/.style={rectangle,draw=zerodds-accent,rounded corners=2pt,
                 minimum width=2.0cm,minimum height=0.7cm,align=center,
                 fill=white},
  arrow/.style={-{Stealth[length=2.0mm]},thick}
]
  % Centre the Endpoint between Client A and Client B so that both
  % west-anchors land on its west edge.
  \node[box, fill=zerodds-primary, text=white,
        minimum height=2.6cm, minimum width=2.6cm]
                                          (endpoint) {DDS-AMQP\\Endpoint\\(Endpoint\\Profile)};
  \node[box, left=2.4cm of endpoint, yshift=0.8cm]   (client1) {AMQP Client A\\(Edge Sensor)};
  \node[box, left=2.4cm of endpoint, yshift=-0.8cm]  (client2) {AMQP Client B\\(Cloud Service)};
  \node[ddsbox, right=of endpoint, yshift=0.55cm]    (rd) {DDS DataReader};
  \node[ddsbox, right=of endpoint, yshift=-0.55cm]   (wr) {DDS DataWriter};
  \node[box, right=of rd]            (action) {Action\\Planner};
  \node[box, right=of wr]            (sub)    {DDS\\Subscriber};
  \draw[arrow] (client1.east) -- node[above, font=\scriptsize]{AMQP} (endpoint.west|-client1);
  \draw[arrow] (endpoint.west|-client2) -- (client2.east);
  \draw[arrow] (endpoint.east|-rd) -- (rd);
  \draw[arrow] (wr) -- (endpoint.east|-wr);
  \draw[arrow] (rd) -- (action);
  \draw[arrow] (sub) -- (wr);
\end{tikzpicture}%
}
\end{center}

\subsection{Sidecar Topology}\label{sec:topology-sidecar}

In the Sidecar topology, the DDS-AMQP bridge is itself an AMQP
client. It connects to an external broker (e.g.\ RabbitMQ, Solace
PubSub+) and translates between the broker and DDS. This topology
preserves the broker's queueing, persistence, and routing
semantics; the bridge replaces a per-application connector layer
that would otherwise be implemented separately.

\begin{center}
\resizebox{\textwidth}{!}{%
\begin{tikzpicture}[
  node distance=0.8cm and 0.7cm,
  every node/.style={font=\small},
  box/.style={rectangle,draw=zerodds-primary,rounded corners=2pt,
              minimum width=2.0cm,minimum height=0.9cm,align=center,
              fill=zerodds-light},
  ddsbox/.style={rectangle,draw=zerodds-accent,rounded corners=2pt,
                 minimum width=2.0cm,minimum height=0.7cm,align=center,
                 fill=white},
  arrow/.style={-{Stealth[length=2.0mm]},thick}
]
  \node[box]                       (producer) {AMQP Producer\\(Edge Device)};
  \node[box, right=of producer]    (broker)   {RabbitMQ /\\Solace};
  \node[box, right=of broker, fill=zerodds-primary, text=white]
                                   (bridge)   {DDS-AMQP\\Bridge\\(Bridge\\Profile)};
  \node[ddsbox, right=of bridge]   (topic)    {DDS Topic};
  \node[box, right=of topic]       (planner)  {Action\\Planner};
  \draw[arrow] (producer) -- (broker);
  \draw[arrow] (broker)   -- (bridge);
  \draw[arrow] (bridge)   -- (topic);
  \draw[arrow] (topic)    -- (planner);
\end{tikzpicture}%
}
\end{center}

\subsection{Hybrid Topology}\label{sec:topology-hybrid}

In the Hybrid topology, an implementation simultaneously realises
both the Endpoint and the Bridge profiles within a single process.
This permits fan-in/fan-out architectures where some AMQP clients
connect directly while others remain connected through an external
broker. The DDS side is unified; AMQP clients of either kind are
served from the same DDS topic catalog.

\begin{center}
\resizebox{\textwidth}{!}{%
\begin{tikzpicture}[
  node distance=0.7cm and 1.0cm,
  every node/.style={font=\small},
  box/.style={rectangle,draw=zerodds-primary,rounded corners=2pt,
              minimum width=2.2cm,minimum height=0.9cm,align=center,
              fill=zerodds-light},
  ddsbox/.style={rectangle,draw=zerodds-accent,rounded corners=2pt,
                 minimum width=2.2cm,minimum height=0.8cm,align=center,
                 fill=white},
  arrow/.style={-{Stealth[length=2.0mm]},thick}
]
  % Anchor on the Hybrid endpoint and place peers relative to it.
  \node[box, fill=zerodds-primary, text=white,
        minimum height=3.0cm, minimum width=2.6cm]  (hybrid)
        {DDS-AMQP\\Endpoint+\\Bridge\\(Hybrid)};
  \node[box, left=of hybrid, yshift=-1.0cm]         (broker) {External\\Broker};
  \node[box, left=of broker]                        (legacy) {Legacy AMQP\\client};
  \node[box]  at ([yshift=1.0cm]hybrid.west -| legacy)  (direct) {Direct AMQP\\client};
  \node[ddsbox, right=of hybrid]                    (topic)  {DDS Topic};
  \draw[arrow] (direct.east) -- (hybrid.west|-direct);
  \draw[arrow] (legacy.east) -- (broker.west);
  \draw[arrow] (broker.east) -- (hybrid.west|-broker);
  \draw[arrow] (hybrid.east) -- (topic);
\end{tikzpicture}%
}
\end{center}

\section{Use-Case Examples}

\subsection{Edge Sensor to Cloud Subscriber}

An edge device (sensor gateway, embedded controller) emits sensor results
via an AMQP-1.0 producer library. A DDS-AMQP endpoint receives
these emissions and writes them as DDS samples on a configured
topic. A downstream DDS subscriber (e.g.\ an action-planning
component) consumes them. No intermediate broker is required.

\subsection{Cloud Service to Edge Subscriber}

A cloud-side AMQP-1.0 producer (e.g.\ Microsoft Azure Service Bus
SDK) publishes commands. A DDS-AMQP endpoint receives them and
writes them onto an edge DDS topic, where edge subscribers consume
them. The same endpoint may also expose other DDS topics as AMQP
sources for cloud-side telemetry consumption.

\subsection{Multi-Vendor Federation}

Two DDS-AMQP endpoints, each connected to a separate DDS domain,
exchange traffic via direct AMQP link attachment. The endpoints
discover each other's exposed topics through the address-catalog
mechanism of~\S\,\ref{sec:pim-discovery-bridging}.

\subsection{Migration from Connector Application}

An existing deployment has an AMQP-only sensor producer, an AMQP
broker, and a custom connector application that consumes AMQP and
publishes to DDS. The Sidecar topology of~\S\,\ref{sec:topology-sidecar}
permits replacing the custom connector with a Bridge-Profile
DDS-AMQP implementation. Neither the sensor nor the broker is
affected.

% ============================================================================
%  Chapter 7 — Platform-Independent Model (PIM).
% ============================================================================
\chapter{Platform-Independent Model (Normative)}\label{ch:pim}

This chapter defines the normative mapping between DDS and AMQP
abstractions, independent of any wire format. The platform-specific
realisation is given in Chapter~\ref{ch:psm}.

\section{Type-System Mapping}\label{sec:pim-type-mapping}

The DDS-AMQP type-system mapping is defined element-by-element from
the OMG IDL primitive types via XCDR2 to the AMQP primitive
type-codes of \omgref{amqp-1.0-types}{\S\,1.6}. The mapping is
total over IDL primitives; unmapped IDL types \rfckw{SHALL} cause
the encoder to raise an \texttt{UnsupportedType} error.

\subsection{Primitive Mapping}

% Use sans-serif (\textsf) for hex codes so '0' and 'O' (and '1'
% and 'l') are unambiguously distinguishable; whitespace via
% \enspace between separator and codes prevents any visual
% concatenation in the rendered PDF.
\begin{longtable}{p{3.0cm}p{1.8cm}p{3.6cm}p{3.2cm}}
  \toprule
  \textbf{IDL Type} & \textbf{XCDR2 Width} & \textbf{AMQP Code} & \textbf{Notes} \\
  \midrule
  \endhead
  \texttt{boolean}              & 1 octet       & \textsf{0x41}\enspace/\enspace\textsf{0x42}                                       & compact form \\
  \texttt{octet}                & 1 octet       & \textsf{0x50}                                                                    & ubyte \\
  \texttt{int8}                 & 1 octet       & \textsf{0x51}                                                                    & signed byte \\
  \texttt{uint8}                & 1 octet       & \textsf{0x50}                                                                    & ubyte \\
  \texttt{short}                & 2 octets      & \textsf{0x61}                                                                    & big-endian \\
  \texttt{unsigned short}       & 2 octets      & \textsf{0x60}                                                                    & big-endian \\
  \texttt{long}                 & 4 octets      & \textsf{0x71}\enspace/\enspace\textsf{0x54}                                       & compact form for $-128 \le n \le 127$ \\
  \texttt{unsigned long}        & 4 octets      & \textsf{0x70}\enspace/\enspace\textsf{0x52}\enspace/\enspace\textsf{0x43}            & compact / zero \\
  \texttt{long long}            & 8 octets      & \textsf{0x81}\enspace/\enspace\textsf{0x55}                                       & compact form for $-128 \le n \le 127$ \\
  \texttt{unsigned long long}   & 8 octets      & \textsf{0x80}\enspace/\enspace\textsf{0x53}\enspace/\enspace\textsf{0x44}            & compact / zero \\
  \texttt{float}                & 4 octets      & \textsf{0x72}                                                                    & IEEE 754 binary32 BE \\
  \texttt{double}               & 8 octets      & \textsf{0x82}                                                                    & IEEE 754 binary64 BE \\
  \texttt{long double}          & 16 octets     & \textsf{0x82}                                                                    & narrowed to binary64; see \S\,\ref{sec:long-double-narrowing} \\
  \texttt{char}                 & 1 octet       & \textsf{0x73}                                                                    & UTF-8 ASCII subset; see \S\,\ref{sec:char-encoding} \\
  \texttt{wchar}                & 4 octets      & \textsf{0x73}                                                                    & XCDR2 wire form is 4-octet UCS-4 per X-Types 1.3 §7.4.3; runtime source representation is platform-defined (UTF-16 or UTF-32) per OMG IDL §7.4.1.4.4 and is transcoded to AMQP UTF-32 codepoint per \S\,\ref{sec:char-encoding-detail} \\
  \texttt{string}               & length+bytes  & \textsf{0xA1}\enspace/\enspace\textsf{0xB1}                                       & UTF-8 verbatim \\
  \texttt{wstring}              & length+wchars & \textsf{0xA1}\enspace/\enspace\textsf{0xB1}                                       & transcoded to UTF-8, see \S\,\ref{sec:wstring-encoding} \\
  \bottomrule
\end{longtable}

\subsection{Long Double Narrowing}\label{sec:long-double-narrowing}

OMG IDL \texttt{long double} is platform-defined as IEEE-754
binary128 or x86-extended-precision (binary80), neither of which
has a corresponding AMQP-1.0 wire-type. The encoder \rfckw{SHALL}
narrow \texttt{long double} values to AMQP \texttt{double}
(IEEE-754 binary64, code \textsf{0x82}) and \rfckw{SHALL} record
this narrowing as a \texttt{decode-error}-class diagnostic
(non-fatal warning) when the narrowed value differs from the
source by more than the binary64 unit-in-the-last-place.

A receiver that requires full \texttt{long double} precision
\rfckw{SHALL} use the AMQP \texttt{binary} mapping (16 octets
verbatim from the source representation) by configuring the
topic's encoding mode to \texttt{MODE\_PASSTHROUGH}; the
binary64-narrowing form is the default for JSON and AMQP-Native
modes which lack a 128-bit float wire-type.

\subsection{Endianness Handling}\label{sec:endianness}

DDS XCDR2 (\omgref{OMG X-Types}{\S\,7.4.3}) supports both
big-endian (BE) and little-endian (LE) sample encodings; the
encoding flag is carried in the encapsulation header at the
start of the serialised sample. AMQP-1.0 numeric types are
strictly big-endian (\omgref{amqp-1.0-types}{\S\,1.4}).

Endpoints \rfckw{SHALL} handle endianness as follows per encoding
mode:

\begin{longtable}{p{3.5cm}p{8.5cm}}
  \toprule
  \textbf{Encoding Mode} & \textbf{Endianness Behaviour} \\
  \midrule
  \endhead
  Pass-Through         & XCDR2 octets are forwarded verbatim into the AMQP \texttt{data} section. The XCDR2 encapsulation header retains the original LE/BE flag, so receivers reconstruct the original encoding. The endpoint \rfckw{SHALL NOT} re-byte-order Pass-Through bytes. \\
  JSON                 & XCDR2 is decoded according to its encapsulation flag, then re-encoded as RFC-8259 JSON (which is byte-order-agnostic). Endianness has no AMQP-side effect. \\
  AMQP-Native          & XCDR2 is decoded according to its encapsulation flag, then re-encoded with each numeric value as a big-endian AMQP-1.0 typed value (per the type-mapping table of \S\,\ref{sec:pim-type-mapping}). The endpoint \rfckw{SHALL} perform the byte-order conversion at the field level; receivers see only AMQP-canonical big-endian values. \\
  \bottomrule
\end{longtable}

\subsection{Character and String Transcoding}\label{sec:char-encoding}

\subsubsection{IDL \texttt{char} (8-bit) to AMQP \texttt{char} (32-bit)}\label{sec:char-encoding-detail}

OMG IDL \texttt{char} is an 8-bit value. To remain consistent with
the \texttt{string} type's UTF-8 wire form
(\S\,\ref{sec:wstring-encoding}), this specification restricts
\texttt{char} to the UTF-8 single-byte subset, i.e.\ ASCII
codepoints \texttt{0x00}--\texttt{0x7F}.

The endpoint \rfckw{SHALL} widen an in-range IDL \texttt{char}
(byte value $\le$ \texttt{0x7F}) to an AMQP \texttt{char} (code
\textsf{0x73}, 4-octet big-endian Unicode codepoint per
\omgref{amqp-1.0-types}{\S\,1.6.16}) by zero-extension: source
byte \texttt{0xNN} produces AMQP codepoint
\texttt{0x0000\,00NN}.

A source byte $>$ \texttt{0x7F} \rfckw{SHALL} cause the encoder
to emit a \texttt{decode-error}-class fault
(\S\,\ref{ch:error-conditions}). Applications that require
non-ASCII single-character data \rfckw{MUST} use IDL
\texttt{wchar}, which carries the full Unicode-codepoint range.

\begin{specnote}
Earlier drafts permitted ISO-8859-1 source encoding for
\texttt{char}, which produced inconsistent behaviour on
non-ASCII content because \texttt{string} is strictly UTF-8.
v1.0.0-beta1 retracts the ISO-8859-1 default and aligns
\texttt{char} with the UTF-8 ASCII subset.
\end{specnote}

\begin{specnote}
\textbf{Migration from Latin-1 IDL codebases.} Existing IDL
codebases that store Latin-1 (ISO-8859-1) bytes in \texttt{char}
fields — typically Western European text with diacritics — will
fail at the encoder boundary under this revision. The supported
migration paths are: (i)~widen the affected fields to
\texttt{wchar} (preserves all codepoints, doubles or quadruples
the field width on the wire); (ii)~widen to \texttt{string}
(preserves all codepoints, variable width); or (iii)~apply an
application-level transcode from Latin-1 to UTF-8 \texttt{string}
on the producer side. Migration via (i) is IDL-source-compatible
on most consumers because \texttt{wchar} accepts the full Latin-1
range as a subset of Unicode BMP. The choice is deployment-local
and is not constrained by this specification.
\end{specnote}

\subsubsection{IDL \texttt{wchar} to AMQP \texttt{char}}

OMG IDL \texttt{wchar} is platform-defined as either UTF-16 or
UTF-32. The endpoint \rfckw{SHALL} interpret \texttt{wchar}
according to the type's declared encoding (carried in the
TypeObject's \texttt{ExtendedAnnotationParameterValue} per
\omgref{OMG X-Types}{\S\,7.2.2.4.6}); when the encoding is
UTF-16 with a surrogate-pair, the codepoint \rfckw{SHALL} be
reconstructed before encoding into the AMQP \texttt{char}.

\subsection{String Transcoding}\label{sec:wstring-encoding}

\subsubsection{IDL \texttt{string} (UTF-8) to AMQP \texttt{string}}

OMG IDL \texttt{string} is UTF-8-encoded by current convention
(see \omgref{OMG IDL}{\S\,7.4.1.4.4.4}). The endpoint
\rfckw{SHALL} pass the octets through verbatim into the AMQP
\texttt{string} (codes \textsf{0xA1}\,/\,\textsf{0xB1}) which
itself is normatively UTF-8
(\omgref{amqp-1.0-types}{\S\,1.6.20}).

\subsubsection{IDL \texttt{wstring} to AMQP \texttt{string} (transcoding)}

OMG IDL \texttt{wstring} is platform-defined as either UTF-16 or
UTF-32. AMQP \texttt{string} is strictly UTF-8. The endpoint
\rfckw{MUST} transcode the wstring to UTF-8 before encoding;
implementations \rfckw{MUST NOT} place raw UTF-16 or UTF-32
octets into the AMQP \texttt{string} body (this would cause
strict AMQP consumers to reject the message as malformed).

The transcoding \rfckw{SHALL} follow the standard Unicode
algorithm:

\begin{enumerate}
  \item For UTF-16 source, decode surrogate pairs to codepoints,
        then encode each codepoint as UTF-8 per RFC 3629.
  \item For UTF-32 source, encode each codepoint directly as
        UTF-8 per RFC 3629.
  \item Reject any source that contains an invalid surrogate or
        a codepoint $\geq$ \texttt{0x11\,0000} with a
        \texttt{decode-error} fault
        (\S\,\ref{ch:error-conditions}).
\end{enumerate}

\subsection{Time Types}\label{sec:pim-time-types}

The DDS \texttt{Time\_t} and \texttt{Duration\_t} carry both a
seconds and a nanoseconds field; AMQP \texttt{timestamp} (code
\texttt{0x83}) is millisecond-resolution. Two mapping forms are
defined; the choice between them is governed by the
\texttt{time\_mapping}-flag in the topic configuration
(Annex~A).

\subsubsection{Vendor-Specific Descriptor Encoding}\label{sec:vendor-descriptor-encoding}

OASIS AMQP 1.0 (\omgref{amqp-1.0-types}{\S\,1.5.2}) requires
that numeric described-composite descriptors use a vendor-id
allocated through the IANA Private Enterprise Number (PEN)
registry. The ZeroDDS Project does not currently hold an IANA
PEN allocation; consequently, this specification \rfckw{SHALL
NOT} use numeric ulong descriptors for ZeroDDS-defined
composites. All ZeroDDS-defined described-composites
\rfckw{SHALL} use \emph{symbolic} descriptors instead
(\omgref{amqp-1.0-types}{\S\,1.5.3}).

\subsubsection{Symbolic Descriptor Convention}

A ZeroDDS-defined described composite \rfckw{SHALL} carry an
AMQP \texttt{symbol} descriptor (code \textsf{0xA3} for short
form $\le$\,255 octets, \textsf{0xB3} for long form) whose ASCII
content matches the following grammar:

\begin{lstlisting}
descriptor   ::= "zerodds:" name [":" version]
name         ::= identifier
version      ::= unsigned-integer
identifier   ::= [A-Za-z][A-Za-z0-9_-]*
\end{lstlisting}

Identifiers are case-sensitive; the namespace prefix
\texttt{"zerodds:"} \rfckw{SHALL} be lower-case.

\subsubsection{Reserved Symbolic Descriptors}

The following symbol descriptors are reserved for use by this
specification:

\begin{longtable}{p{4.5cm}p{8cm}}
  \toprule
  \textbf{Descriptor symbol} & \textbf{Meaning} \\
  \midrule
  \endhead
  \texttt{"zerodds:topic-mapping:1"} & Address-Catalog topic-mapping entry (\S\,\ref{sec:pim-discovery-bridging}) \\
  \texttt{"zerodds:time:1"}          & High-resolution time composite (\S\,\ref{sec:pim-time-types}, opt-in) \\
  \texttt{"zerodds:dispose:1"}       & Reserved for future revisions \\
  \bottomrule
\end{longtable}

The trailing \texttt{:1} is the version suffix; a future
revision that changes the body layout of a composite \rfckw{SHALL}
increment the version (e.g.\ \texttt{"zerodds:time:2"}). A
receiver that does not recognise the version \rfckw{SHALL} treat
the composite as opaque and reject the transfer with
\texttt{amqp:not-implemented} if normative semantics are
required.

\begin{specnote}
Earlier drafts of v1.0.0-beta1 used numeric ulong descriptors of
the form \texttt{0x0000\,006c\,XXXX\,XXXX}. PEN \texttt{0x6c}
(decimal 108) is not allocated to the ZeroDDS Project. Symbolic
descriptors avoid the IANA PEN coordination requirement entirely
and are equally well-supported by AMQP-1.0 libraries.
Implementations of earlier drafts \rfckw{SHALL} migrate to
symbolic descriptors before claiming v1.0.0-beta1 conformance.
\end{specnote}

\subsubsection{Standard Mapping (default, normative)}

By default, DDS \texttt{Time\_t} \rfckw{SHALL} be mapped to the
standard AMQP \texttt{timestamp} (code \texttt{0x83}, 8-byte BE
signed milliseconds since Unix epoch). The DDS sample's
\texttt{Time\_t.sec*1000 + Time\_t.nanosec/1\,000\,000}
\rfckw{SHALL} be placed into the Properties section's
\texttt{creation-time} field. The nanoseconds remainder
($\texttt{Time\_t.nanosec} \bmod 1\,000\,000$) \rfckw{SHALL} be
recorded in the Application-Property \texttt{dds:nsec}
(see~\S\,\ref{sec:psm-app-properties}) when sub-millisecond
precision is required by the application.

This form is interoperable with any AMQP-1.0 client; AMQP-aware
brokers that order or filter by \texttt{creation-time} work
unchanged. Applications that need sub-millisecond precision read
\texttt{dds:nsec} after \texttt{creation-time}.

\subsubsection{High-Resolution Composite Mapping (opt-in, normative)}

When the topic's configuration sets
\texttt{time\_mapping = MAPPING\_COMPOSITE} (Annex~A), DDS
\texttt{Time\_t} \rfckw{SHALL} be mapped to an AMQP described
composite with symbolic descriptor \texttt{"zerodds:time:1"}
(see~\S\,\ref{sec:vendor-descriptor-encoding}) whose body is a
2-element list of (\texttt{long} seconds, \texttt{uint}
nanoseconds). This form preserves full nanosecond precision and
is intended for DDS-to-DDS federation through AMQP where both
sides are DDS-AMQP-aware. Generic AMQP consumers \rfckw{SHALL
NOT} be required to decode the composite.

\begin{specnote}
v1.0.0-beta1 inverts the precedence of the two forms relative to
earlier drafts. The standard mapping is the new default because
it preserves interoperability with non-DDS-aware AMQP consumers.
The composite form remains available as opt-in when both ends are
known to be DDS-AMQP-aware.
\end{specnote}

\subsection{Identifier Types}\label{sec:pim-identifier-types}

\subsubsection{InstanceHandle (DDS-opaque)}

The DDS \texttt{InstanceHandle\_t} is a 16-octet opaque value that
is \emph{not} guaranteed to satisfy the RFC-4122 UUID
version/variant bit-layout. To remain interoperable with strict
AMQP consumers that validate UUID bits, implementations
\rfckw{SHALL} map \texttt{InstanceHandle\_t} to an AMQP
\texttt{binary} of length 16 (code \textsf{0xA0} when length
$\le$ 255 --- always the short form for 16-byte values).
Implementations \rfckw{SHALL NOT} use AMQP \texttt{uuid}
(code \textsf{0x98}) for \texttt{InstanceHandle\_t}.

\subsubsection{Security GUID (RFC-4122 conformant)}

The DDS-Security \texttt{GUID\_t} as defined in
\omgref{OMG DDS-Security}{\S\,9.4.2.3} is constructed to satisfy
RFC~4122 version/variant bits. \texttt{GUID\_t} \rfckw{SHALL} be
mapped to AMQP \texttt{uuid} (code \textsf{0x98}, 16 bytes).

\subsubsection{Rationale}

Splitting the mapping by RFC-4122 conformance permits strict AMQP
consumers (e.g.\ those using a UUID-validation library) to accept
DDS-Security identities while treating
\texttt{InstanceHandle\_t} as opaque binary. Earlier drafts of
this specification mapped both to \texttt{uuid}; that mapping has
been retracted as of v1.0.0-beta1 because it caused interoperability
failures with strict consumers.

\subsection{Sequence and Array}

A DDS \texttt{sequence$\langle$T$\rangle$} or fixed-size array of
\texttt{T} \rfckw{SHALL} be mapped to either an AMQP \texttt{list}
(code \texttt{0xC0/0xD0}) or an AMQP \texttt{array} (code
\texttt{0xE0/0xF0}). The \texttt{array} form is \rfckw{RECOMMENDED}
when all elements share a single AMQP primitive type-code, as it
is more compact on the wire; the \texttt{list} form is
\rfckw{REQUIRED} when elements differ in type or are themselves
described composites. An empty sequence or array \rfckw{SHALL}
be encoded as the \texttt{list0} type-code \texttt{0x45} (single
octet, no size and no count fields), \emph{not} as an empty
\texttt{0xC0} or \texttt{0xE0} list. Decoders \rfckw{SHALL}
accept \texttt{0x45} as semantically equivalent to a zero-element
\texttt{0xC0}.

\subsection{Map}

A DDS map type (X-Types \texttt{TK\_MAP}) \rfckw{SHALL} be mapped
to an AMQP \texttt{map} (code \textsf{0xC1}\,/\,\textsf{0xD1}).

OASIS AMQP 1.0 \omgref{amqp-1.0-types}{\S\,1.6.23}
explicitly states that \emph{the order of distinct keys is not
significant in the encoding}; conformant AMQP libraries
\rfckw{MAY} reorder map entries during parse and re-serialise.
Applications that depend on key-order preservation (e.g.\ DDS
maps treated as ordered associative arrays) \rfckw{MUST NOT} rely
on the AMQP \texttt{map} encoding and \rfckw{SHALL} instead
encode such data as an AMQP \texttt{list} of two-element lists
\texttt{(key, value)}, where outer-list order is preserved per
\omgref{amqp-1.0-types}{\S\,1.6.22}. The choice of encoding
\rfckw{SHALL} be made by the type author at IDL design time and
recorded in the type's \texttt{TypeObject}.

\begin{specnote}
The previous draft required \texttt{insertion order} preservation
in the \texttt{map} encoding. That requirement is unfulfillable
with standard AMQP tools; v1.0.0-beta1 retracts it and provides
the list-of-pairs alternative for ordered-map use-cases.
\end{specnote}

\section{Composite-Type Mapping}\label{sec:pim-composite-mapping}

A DDS composite type (\texttt{struct}, \texttt{union}, \texttt{valuetype})
\rfckw{SHALL} be mapped to an AMQP described composite consisting of
a descriptor and a body.

\subsection{Descriptor}

The descriptor of a DDS composite \rfckw{SHALL} be derived from
the X-Types \texttt{TypeIdentifier} of the type
(\omgref{OMG X-Types}{\S\,7.3.4}). Two encoding forms are
defined; implementations \rfckw{SHALL} support both and select
based on the configuration setting
\texttt{descriptor\_form} (Annex~A).

\subsubsection{Truncated Form (default)}

The first eight octets of the type's equivalence-hash
\rfckw{SHALL} be interpreted as a big-endian unsigned 64-bit
integer and used as an AMQP \texttt{ulong} descriptor (code
\textsf{0x80}). This form is space-efficient and matches the
described-composite convention used by the AMQP specification
itself for performatives.

\subsubsection{Full-Identifier Form (collision-safe)}

The full 14-octet \texttt{TypeIdentifier} (or longer for
\texttt{TypeIdentifier} variants that include hash extensions)
\rfckw{SHALL} be encoded as an AMQP \texttt{symbol} descriptor
whose value is the lower-case hexadecimal serialisation of the
identifier bytes, prefixed with \texttt{"dds:type:"}.

\begin{specexample}
A type whose 14-byte TypeIdentifier is
\texttt{0xDEADBEEF\ldots} produces the symbol
\texttt{dds:type:deadbeef\ldots}.
\end{specexample}

\subsubsection{Collision Risk and Mitigation}

The truncated 64-bit form has a birthday-bound collision
probability of approximately $2^{-32}$ samples
(\rfckw{ROUGHLY} 4~billion distinct types). Deployments that
register more than \rfckw{ROUGHLY} $10^9$ distinct topic types
\rfckw{SHOULD} configure \texttt{descriptor\_form = DESC\_FULL}.

When \texttt{descriptor\_form = DESC\_TRUNCATED}, the sending
endpoint \rfckw{SHALL} additionally populate the
Application-Property \texttt{dds:type-id}
(\S\,\ref{sec:psm-app-properties}) with the full 14-byte
TypeIdentifier hex string, and the receiving endpoint
\rfckw{SHALL} inspect it. When the descriptor matches a
locally-known type but \texttt{dds:type-id} differs, the
endpoint \rfckw{SHALL} reject the transfer with disposition
\texttt{rejected} and error condition
\texttt{amqp:decode-error}. When \texttt{descriptor\_form =
DESC\_FULL} (the 14-byte symbolic form), \texttt{dds:type-id}
\rfckw{MAY} be omitted because the descriptor itself already
carries the full identifier.

\begin{specnote}
This descriptor convention establishes an isomorphism between
X-Types-aware DDS endpoints and AMQP-typed consumers. AMQP-only
consumers without X-Types support \rfckw{MAY} treat the descriptor
as opaque.
\end{specnote}

\subsection{Body}

The body of a DDS composite \rfckw{SHALL} be encoded as an AMQP
\texttt{list} (code \textsf{0xC0}\,/\,\textsf{0xD0}) whose
elements correspond, in IDL declaration order, to the fields of
the composite type.

\subsubsection{Absent Optional Fields}

A field annotated \texttt{@optional} whose runtime value is absent
\rfckw{SHALL} be encoded as the AMQP \texttt{null} value (code
\textsf{0x40}) at its positional slot in the list. The field
\rfckw{MUST NOT} be omitted from the list, because positional
correspondence by IDL declaration order is the only addressing
mechanism into the body.

\subsubsection{Trailing Absent Fields}

A run of trailing absent optional fields \rfckw{MAY} be encoded by
shortening the list (omitting trailing \texttt{null} values).
Decoders \rfckw{SHALL} treat any positions beyond the encoded
list-length as \texttt{null}. This optimisation matches AMQP's
list-of-mixed-types convention and is used by AMQP itself for
encoding performatives with optional trailing fields.

\subsubsection{Nested Composites}

Nested composite types (struct fields whose type is itself a
composite) \rfckw{SHALL} themselves be encoded as described
composites recursively per \S\,\ref{sec:pim-composite-mapping}.

\subsubsection{Mutable Types}

For an \texttt{@mutable} composite type, the body \rfckw{SHALL}
be an AMQP \texttt{map} (code \textsf{0xC1}\,/\,\textsf{0xD1})
whose keys are the integer member-IDs declared by
\texttt{@id(N)} annotations and whose values are the encoded
field values. The list-positional form is invalid for
\texttt{@mutable} types because their member-set evolves across
versions.

\subsection{Union}

A DDS \texttt{union} \rfckw{SHALL} be encoded as a described
composite whose body is a 2-element list: the discriminator value
and the active branch's value (or \texttt{null} if the branch is
empty).

\section{Address-Resolution Model}\label{sec:pim-address-resolution}

An AMQP source address (for an outgoing-from-endpoint link) or
target address (for an incoming-to-endpoint link)
\rfckw{SHALL} be mapped to a DDS Topic name in accordance with the
following rules. The address grammar is the primary mechanism for
multi-domain routing: the \texttt{domain://}\,$N$\,/\,$\textit{topic}$
form selects DDS domain $N$ explicitly, allowing a single
endpoint to expose topics from any number of DDS domains
simultaneously without ambiguity in the AMQP namespace.

\begin{enumerate}
  \item A bare address (e.g.\ \texttt{tracking}) \rfckw{SHALL} resolve
        to the DDS Topic of the same name in the endpoint's default
        domain, with the empty Partition.
  \item An address of the form
        \texttt{domain://}$\langle$\texttt{n}$\rangle$\texttt{/}$\langle$\texttt{topic}$\rangle$
        \rfckw{SHALL} resolve to the named topic in DDS domain $n$
        ($n \in [0, 232]$).
  \item An address of the form
        \texttt{domain://}$\langle$\texttt{n}$\rangle$\texttt{/}$\langle$\texttt{topic}$\rangle$\texttt{?partition=}$\langle$\texttt{p}$\rangle$
        \rfckw{SHALL} additionally constrain the Partition QoS of
        the resulting Publisher/Subscriber to value $p$.
  \item An address containing a wildcard character (\texttt{*} or
        \texttt{?}) \rfckw{SHALL} be interpreted as a Partition
        pattern as defined in \omgref{OMG DDS}{\S\,2.2.3.21}.
\end{enumerate}

The endpoint configuration \rfckw{MAY} additionally define static
address aliases (see Annex~A). Static aliases override rules~1--4
when the address matches an alias name.

\section{QoS-Policy Mapping}\label{sec:pim-settlement}

This section defines the mapping of every DDS QoS policy that is
observable at the AMQP boundary. Policies not listed in this
section have no AMQP-side semantics and \rfckw{SHALL} be honoured
by the DDS-side endpoint without any AMQP-visible effect.

\subsection{QoS in the Absence of Broker Functionality}

DDS-AMQP explicitly excludes broker-class features (queue
persistence, exchange routing, dead-letter, message-TTL) from its
scope (\S\,\ref{sec:out-of-scope}). Several DDS QoS policies
nevertheless have semantic content that requires durable,
broker-mediated machinery to fully translate into a multi-hop
AMQP path. The following table records, for each affected QoS,
how the DDS guarantee is preserved when no broker is present:

\begin{longtable}{p{3.5cm}p{8.5cm}}
  \toprule
  \textbf{DDS QoS} & \textbf{Behaviour without intermediate broker} \\
  \midrule
  \endhead
  \texttt{LIFESPAN}        & The endpoint enforces LIFESPAN locally on the DDS side; the AMQP \texttt{absolute-expiry-time} is propagated, but if the AMQP path traverses no broker, expiry is the receiving endpoint's responsibility. \\
  \texttt{TRANSIENT\_LOCAL DURABILITY} & The endpoint replays its own DataReader's transient-local cache to a newly-attached AMQP receiver; this is local to the endpoint and survives no endpoint restart. \\
  \texttt{TRANSIENT} / \texttt{PERSISTENT DURABILITY} & Cannot be honoured without broker functionality; rejected with \texttt{amqp:not-implemented} as defined in \S\,\ref{sec:pim-settlement}. \\
  \texttt{DEADLINE}        & DDS-side deadline-tracking is local to the endpoint; missed deadlines are signalled via the catalog channel (\S\,\ref{sec:pim-discovery-bridging}) but cannot be retro-actively enforced once the sample has crossed into AMQP. \\
  \texttt{HISTORY}         & DCPS per-instance cache; entirely DDS-side; AMQP receives only the materialised stream of samples that DDS itself emits per its HISTORY-policy. \\
  \bottomrule
\end{longtable}

Deployments that require any of the broker-class guarantees
\rfckw{SHOULD} use the Sidecar topology
(\S\,\ref{sec:topology-sidecar}) with a broker (e.g.\ RabbitMQ
configured for durability) sitting between the AMQP-clients and
the DDS-AMQP bridge. The broker then provides the broker-class
features that this specification does not.

\subsection{RELIABILITY}

The mapping is bidirectional and \rfckw{SHALL} hold for both the
DDS-derived link-attribute production direction (endpoint
constructs the \texttt{attach} performative based on its
DataWriter/DataReader QoS) and the AMQP-derived QoS-resolution
direction (endpoint chooses DataWriter/DataReader QoS based on
the \texttt{attach} performative received from the peer).

\subsubsection{DDS-to-AMQP (link establishment)}

The endpoint \rfckw{SHALL} set the AMQP link attributes from the
DDS-side QoS as follows. All four AMQP attributes \rfckw{SHALL}
be set on the \texttt{attach} performative when establishing the
link.

\begin{longtable}{p{3.0cm}p{2.4cm}p{2.6cm}p{4.4cm}}
  \toprule
  \textbf{DDS RELIABILITY} & \textbf{snd-settle-mode} & \textbf{rcv-settle-mode} & \textbf{Per-Transfer Behaviour} \\
  \midrule
  \endhead
  \texttt{BEST\_EFFORT} & \texttt{settled} (1) & \texttt{first} (0)
    & sender pre-settles each transfer; receiver \rfckw{MUST NOT} emit a \texttt{disposition}. \\
  \texttt{RELIABLE} & \texttt{unsettled} (0) & \texttt{second} (1)
    & sender leaves each transfer unsettled; receiver \rfckw{MUST} emit a \texttt{disposition} with terminal state (\texttt{accepted} / \texttt{rejected} / \texttt{released}); the sender then \rfckw{MUST} settle the transfer (per \omgref{amqp-1.0-transport}{\S\,2.7.5} \texttt{rcv-settle-mode = second} semantics). \\
  \bottomrule
\end{longtable}

The numeric values in parentheses are the AMQP-1.0
sender-/receiver-settle-mode values per
\omgref{amqp-1.0-transport}{\S\,2.8.1} and
\S\,2.8.2.

\subsubsection{AMQP-to-DDS (QoS resolution from inbound \texttt{attach})}

When a peer attaches a link to the endpoint, the endpoint
\rfckw{SHALL} derive the DDS-side DataWriter/DataReader
RELIABILITY QoS from the \texttt{attach} performative as
follows:

\begin{longtable}{p{3.0cm}p{2.4cm}p{6.6cm}}
  \toprule
  \textbf{snd-settle-mode} & \textbf{rcv-settle-mode} & \textbf{Resulting DDS RELIABILITY} \\
  \midrule
  \endhead
  \texttt{settled} (1) & any & \texttt{BEST\_EFFORT} \\
  \texttt{unsettled} (0) & \texttt{first} (0) & \texttt{BEST\_EFFORT} (peer is willing to accept loss; rcv-settle-mode=first means receiver settles immediately on receipt) \\
  \texttt{unsettled} (0) & \texttt{second} (1) & \texttt{RELIABLE} \\
  \texttt{mixed} (2) & \texttt{second} (1) & \texttt{RELIABLE} (sender chooses per-transfer; receiver's terminal-state-second is the binding constraint) \\
  \texttt{mixed} (2) & \texttt{first} (0) & \texttt{BEST\_EFFORT} \\
  \bottomrule
\end{longtable}

\begin{specnote}
The resolution favours \texttt{BEST\_EFFORT} when the peer's
intent is ambiguous; this matches DDS conservative-defaulting
behaviour (\omgref{OMG DDS}{\S\,2.2.3.13}).
\end{specnote}

\subsection{DURABILITY}

The mapping of DDS DURABILITY to AMQP-1.0
\texttt{source.durable} / \texttt{target.durable} attributes is
normatively specified by the following table.
\texttt{source.durable} is set on a Sender-side terminus
(endpoint exposes DDS samples to AMQP); \texttt{target.durable}
is set on a Receiver-side terminus (endpoint accepts AMQP
transfers and writes to DDS).

\begin{longtable}{p{3.5cm}p{3cm}p{6cm}}
  \toprule
  \textbf{DDS DURABILITY} & \textbf{durable field} & \textbf{Endpoint Behaviour} \\
  \midrule
  \endhead
  \texttt{VOLATILE}         & \texttt{none} (0)
    & no replay; transfers from a sender link \rfckw{SHALL NOT} be redelivered after link detach. \\
  \texttt{TRANSIENT\_LOCAL} & \texttt{configuration} (1)
    & endpoint replays its DataReader's transient-local sample cache to a newly-attached receiver before delivering live samples. \\
  \texttt{TRANSIENT}        & \texttt{unsettled-state} (2)
    & out of scope --- requires broker functionality. Implementations \rfckw{SHALL} reject the link attachment with error \texttt{amqp:not-implemented}. \\
  \texttt{PERSISTENT}       & \texttt{unsettled-state} (2)
    & same rejection rule as \texttt{TRANSIENT}. \\
  \bottomrule
\end{longtable}

The numeric values in parentheses are the AMQP-1.0
\texttt{terminus-durability} values per
\omgref{amqp-1.0-transport}{\S\,3.5.4}.

\subsubsection{AMQP-to-DDS (Durability resolution)}

The endpoint \rfckw{SHALL} derive the DDS-side
DataWriter/DataReader DURABILITY QoS from the inbound
\texttt{attach}:

\begin{longtable}{p{4cm}p{8cm}}
  \toprule
  \textbf{terminus.durable} & \textbf{Resulting DDS DURABILITY} \\
  \midrule
  \endhead
  \texttt{none} (0)              & \texttt{VOLATILE} \\
  \texttt{configuration} (1)     & \texttt{TRANSIENT\_LOCAL} \\
  \texttt{unsettled-state} (2)   & rejected with \texttt{amqp:not-implemented}; this profile does not implement broker-level durability. \\
  \bottomrule
\end{longtable}

\subsection{HISTORY}

DDS HISTORY is a per-instance sample-retention policy on the DDS
side (DCPS \omgref{OMG DDS}{\S\,2.2.3.18}); it is not a transport-
level flow-control mechanism. The endpoint \rfckw{SHALL} honour
HISTORY on the DDS-side DataReader/DataWriter exactly as
configured.

The AMQP-side flow-control behaviour is governed by the AMQP-1.0
flow-credit model (\omgref{amqp-1.0-transport}{\S\,2.6.7}) and is
\emph{decoupled} from HISTORY. Implementations \rfckw{SHALL NOT}
expose HISTORY semantics through AMQP-side mechanisms. The
DDS-side cache discipline (last-N, keep-all) is a black-box to
AMQP consumers.

\begin{specnote}
Earlier drafts conflated DDS HISTORY with AMQP unsettled-delivery
windowing; this was incorrect because HISTORY operates on a
per-instance basis at the DCPS layer, not on transport-level
deliveries. v1.0.0-beta1 corrects this: HISTORY has no AMQP-side
visible effect.
\end{specnote}

\subsection{DEADLINE}

A DDS DataWriter's DEADLINE QoS \rfckw{SHALL} be honoured on the
DDS side. Missed deadlines \rfckw{MAY} be reported across the AMQP
boundary as an unsolicited transfer on the catalog link
(\S\,\ref{sec:pim-discovery-bridging}) carrying a
\texttt{deadline-missed} application-property. AMQP consumers
\rfckw{MUST NOT} rely on the DDS-side deadline-tracking; the AMQP
layer has no equivalent guarantee.

\subsection{LATENCY\_BUDGET}

LATENCY\_BUDGET is informative on the DDS side; the AMQP layer
\rfckw{SHALL} ignore it. Implementations \rfckw{MAY} expose the
configured budget value as an Application-Property
\texttt{dds:latency-budget-ms} (uint32) for diagnostic purposes.

\subsection{OWNERSHIP}

\begin{longtable}{p{3.5cm}p{9cm}}
  \toprule
  \textbf{DDS OWNERSHIP} & \textbf{AMQP Behaviour} \\
  \midrule
  \endhead
  \texttt{SHARED}     & default; no AMQP-side semantics. All bound writers contribute to the AMQP-side stream. \\
  \texttt{EXCLUSIVE}  & only the writer with the highest OWNERSHIP\_STRENGTH at any instant produces transfers; transfers from others \rfckw{SHALL} be silently dropped at the endpoint. \\
  \bottomrule
\end{longtable}

\subsection{LIFESPAN}

A DDS sample's LIFESPAN \rfckw{SHALL} be propagated as the AMQP
\texttt{absolute-expiry-time} property. AMQP brokers
\rfckw{SHOULD} drop expired transfers; this is broker behaviour
governed by AMQP and not by this specification.

\subsection{PARTITION}

DDS PARTITION QoS is a \texttt{sequence<string>}: a Publisher or
Subscriber may belong to multiple named partitions
simultaneously. PARTITION \rfckw{SHALL} be honoured by the
address-resolution mechanism
(\S\,\ref{sec:pim-address-resolution}).

A multi-valued PARTITION \rfckw{SHALL} be expressed as repeated
\texttt{?partition=} query parameters in the AMQP address URI
(\texttt{?partition=p1\&partition=p2\&\ldots}). Each occurrence
contributes one element to the bound partition sequence in the
order given. Partition values \rfckw{SHALL} be percent-encoded
per RFC 3986 if they contain reserved or non-ASCII characters.

A single-valued PARTITION \rfckw{MAY} use the equivalent
single-occurrence form \texttt{?partition=p}.

\subsection{All Other QoS Policies}

DDS QoS policies not listed in \S\,\ref{sec:pim-settlement} are
\emph{transparent} to the AMQP boundary: the DDS-side endpoint
\rfckw{SHALL} honour them as configured, but they have no
AMQP-visible behaviour. This includes (non-exhaustive):
\texttt{TIME\_BASED\_FILTER}, \texttt{TRANSPORT\_PRIORITY},
\texttt{LIVELINESS}, \texttt{DESTINATION\_ORDER},
\texttt{RESOURCE\_LIMITS}, \texttt{ENTITY\_FACTORY},
\texttt{WRITER\_DATA\_LIFECYCLE}, \texttt{READER\_DATA\_LIFECYCLE},
\texttt{USER\_DATA}, \texttt{TOPIC\_DATA}, \texttt{GROUP\_DATA},
\texttt{PRESENTATION}, \texttt{TYPE\_CONSISTENCY\_ENFORCEMENT}.

A future revision of this specification \rfckw{MAY} introduce
explicit AMQP-side mappings for any of these policies. Until then
implementations \rfckw{SHALL NOT} expose them via AMQP-specific
properties without prefixing them with \texttt{dds:} (see
\S\,\ref{sec:psm-app-properties}).

\section{Discovery Bridging}\label{sec:pim-discovery-bridging}

A DDS-AMQP endpoint \rfckw{SHALL} expose its set of mapped topics as
an Address Catalog. The catalog \rfckw{SHALL} be retrievable by an
AMQP client by attaching a Receiver link to the source address
\texttt{\$catalog}.

The catalog \rfckw{SHALL} be encoded as a sequence of AMQP
described composites with symbolic descriptor
\texttt{"zerodds:topic-mapping:1"}
(see~\S\,\ref{sec:vendor-descriptor-encoding}). Each entry
\rfckw{SHALL} contain:

\begin{itemize}
  \item AMQP address (string)
  \item DDS topic name (string)
  \item DDS type name (string)
  \item TypeIdentifier — encoded as AMQP \texttt{symbol} (code \textsf{0xA3}\,/\,\textsf{0xB3}) when the topic's \texttt{descriptor\_form = DESC\_FULL} (full 14-octet symbolic identifier), or as AMQP \texttt{ulong} (code \textsf{0x80}, 8-octet big-endian unsigned 64-bit integer) when \texttt{descriptor\_form = DESC\_TRUNCATED}. The \texttt{ulong} value \rfckw{SHALL} be byte-identical to the descriptor used for the same type in~\S\,\ref{sec:pim-composite-mapping}; this is a strict AMQP-type-level consistency, not merely a byte-level one. The catalog entry \rfckw{SHALL NOT} encode a TypeIdentifier as \texttt{uuid}.
  \item Reachable directions: \texttt{producer-to-dds} /
        \texttt{dds-to-consumer} / \texttt{both}
\end{itemize}

The endpoint \rfckw{SHOULD} additionally announce changes to the
catalog by sending unsolicited transfers on the \texttt{\$catalog}
link to attached subscribers. The catalog is the AMQP-visible
counterpart of the DDS SPDP/SEDP mechanism for topic discovery.

\subsection{Behaviour for Unknown Addresses}\label{sec:pim-unknown-address}

When an AMQP client attempts to attach a Sender or Receiver link
to an address that does not appear in the catalog, the endpoint
behaviour is determined by the
\texttt{permit\_dynamic\_topics}-flag of the static configuration
(Annex~A):

\begin{longtable}{p{3.5cm}p{9cm}}
  \toprule
  \textbf{Configuration} & \textbf{Endpoint Behaviour} \\
  \midrule
  \endhead
  \texttt{false} (default) & The endpoint \rfckw{SHALL} reject the link attachment with link-error condition \texttt{amqp:not-found} and a description string that names the unknown address. The connection \rfckw{SHALL} remain open for further attachments to known addresses. \\
  \texttt{true}            & The endpoint \rfckw{SHALL} create a new DDS Topic on-the-fly using the address as topic name and a Caller-supplied default type registered in the configuration (\texttt{dynamic\_topic\_default\_type}). If no default type is configured, behaviour falls back to the \texttt{false} case. The newly created topic \rfckw{SHALL} be appended to the catalog and announced via unsolicited \texttt{\$catalog}-Transfer to all attached catalog subscribers. \\
  \bottomrule
\end{longtable}

\begin{specnote}
Default-rejection (\texttt{permit\_dynamic\_topics = false}) is
chosen as the default to keep the deployment surface predictable.
Operators who wish to permit dynamic topic creation explicitly
configure the flag and accept the operational consequences (no
schema validation against unknown addresses).
\end{specnote}

\section{Key and Instance Mapping}\label{sec:pim-key-instance}

DDS is a data-centric system: data flowing on a topic is
partitioned into \emph{instances} identified by the topic-type's
\texttt{@key} fields (\omgref{OMG DDS}{\S\,2.2.2.2.5}). The
AMQP-side mapping \rfckw{SHALL} preserve the instance identity so
that DDS subscribers downstream of the endpoint see a consistent
per-instance sample stream.

\subsection{AMQP \texttt{group-id} as Key Carrier (normative)}

For every transferred sample whose topic type declares one or
more \texttt{@key} fields, the endpoint \rfckw{SHALL} derive a
canonical key string from the key-field values and place it in
the AMQP Properties section's \texttt{group-id} field
(\omgref{amqp-1.0-messaging}{\S\,3.2.4}).

The canonical key string \rfckw{SHALL} be the lowercase
hexadecimal encoding of the SHA-256 digest computed over the
XCDR2 KeyHash encapsulation of the sample (X-Types 1.3
\S\,7.6.6.4): the encoder \rfckw{SHALL} produce the
KeyHash-canonical XCDR2 byte stream of the key fields, compute
SHA-256 over those bytes, and render the 32-byte digest as 64
lowercase hex characters. The result is opaque to AMQP routing
but stable, byte-for-byte deterministic across implementations,
and robust against arbitrary key-field content (strings
containing the SOLIDUS character, embedded NUL, etc.).

\begin{specexample}
For an IDL type
\begin{lstlisting}
@nested struct Pos { @key long x; @key long y; };
struct Track { @key long sensor_id; @key Pos position; string label; };
\end{lstlisting}
a sample with \texttt{sensor\_id=7, position=\{x=10, y=20\}}
yields a fixed 64-character hex string of the form
\texttt{group-id = "a3f1...e29c"} (the exact value is determined
by the SHA-256 of the XCDR2 KeyHash bytes for those values).
\end{specexample}

\subsection{Subscriber-Side Reconstruction}

The canonical key string is a one-way hash; an AMQP-1.0
subscriber \rfckw{SHALL NOT} attempt to reconstruct individual
key-field values from \texttt{group-id}. Subscribers \rfckw{MAY}
use \texttt{group-id} as an opaque routing-key for AMQP-side
group-affinity, deduplication, or sharding, and \rfckw{MAY}
correlate it with a separately-published key-mapping channel
(implementation-defined; out of scope for this revision).
Subscribers that need structured access to key-field values
\rfckw{SHALL} obtain them from the sample body (XCDR2-encoded
in Pass-Through mode, JSON in JSON mode, or described composite
in AMQP-Native mode).

\subsection{Topic without Keys}

For an unkeyed topic (no \texttt{@key} fields), the
\texttt{group-id} field \rfckw{SHALL} be omitted. Subscribers
\rfckw{MAY} treat the absence of \texttt{group-id} as the
all-instances signal.

\section{Instance Lifecycle Mapping}\label{sec:pim-instance-lifecycle}

DDS supports explicit instance-lifecycle operations
(\omgref{OMG DDS}{\S\,2.2.2.4.2.4}): \texttt{register\_instance},
\texttt{unregister\_instance}, \texttt{dispose}. The AMQP-side
\rfckw{SHALL} provide normative signal forms for each.

\subsection{Outbound Instance State}

When the endpoint forwards a DDS sample to AMQP, the sample's
\texttt{InstanceState} \rfckw{SHALL} be reflected in the
Application-Property \texttt{dds:instance-state}
(\S\,\ref{sec:psm-app-properties}) using the canonical values
\texttt{alive} / \texttt{not-alive-disposed} /
\texttt{not-alive-no-writers}.

A sample that signals a state-change-only event (no payload)
\rfckw{SHALL} carry a body Data section of length 0 (zero-length
\texttt{vbin8}). Subscribers that observe a zero-length body
\rfckw{SHALL} consult \texttt{dds:instance-state} for the new
state.

\subsection{Inbound Operation Signals}

An AMQP producer \rfckw{MAY} actively trigger DDS instance-
lifecycle operations by setting the Application-Property
\texttt{dds:operation} to one of the following values:

\begin{longtable}{p{3cm}p{9cm}}
  \toprule
  \textbf{Value} & \textbf{Endpoint Behaviour} \\
  \midrule
  \endhead
  \texttt{write}              & default; sample is written. The body \rfckw{MUST} contain the encoded sample. \\
  \texttt{register}           & endpoint calls \texttt{DataWriter::register\_instance} for the key derived from \texttt{group-id} (\S\,\ref{sec:pim-key-instance}). The body \rfckw{MAY} be a partial sample carrying only the key fields; non-key fields \rfckw{SHALL} be ignored. \\
  \texttt{unregister}         & endpoint calls \texttt{DataWriter::unregister\_instance}. The body \rfckw{SHALL} be empty. \\
  \texttt{dispose}            & endpoint calls \texttt{DataWriter::dispose}. The body \rfckw{SHALL} be empty. \\
  \bottomrule
\end{longtable}

The endpoint \rfckw{SHALL} reject any other value for
\texttt{dds:operation} with link error
\texttt{amqp:not-implemented}.

\subsection{Disposition Mapping}

When DDS-side disposal succeeds, the endpoint \rfckw{SHALL}
acknowledge the AMQP transfer with \texttt{accepted} disposition.
Failure (e.g.\ unknown instance for \texttt{unregister}) \rfckw{SHALL}
be acknowledged with \texttt{rejected} disposition carrying error
condition \texttt{amqp:precondition-failed} and a description
string identifying the failed operation.

\section{Conflict Resolution between URI and Properties}\label{sec:pim-conflict-resolution}

Two information channels redundantly carry partition information:
the URI-query syntax \texttt{?partition=} of the address
(\S\,\ref{sec:pim-address-resolution}) and the
Application-Property \texttt{dds:partition}
(\S\,\ref{sec:psm-app-properties}). The following precedence
\rfckw{SHALL} be applied:

\begin{enumerate}
  \item If both the URI-query partition and the
        \texttt{dds:partition} property are present and \emph{equal},
        the partition is honoured.
  \item If both are present and \emph{differ}, the URI-query
        partition \rfckw{SHALL} take precedence as the
        authoritative value, and the Application-Property
        \rfckw{SHALL} be silently overwritten on the DDS side.
        Implementations \rfckw{MAY} log a warning-level
        diagnostic.
  \item If only one of the two is present, that value
        \rfckw{SHALL} be used.
  \item If neither is present, the empty (default) partition
        \rfckw{SHALL} be used.
\end{enumerate}

\begin{specnote}
URI-precedence is chosen because the URI is part of the link
attachment and is therefore more visible to operators; the
Application-Property is per-message and easier for an
application to set inadvertently. Operators typically build
allow/deny rules against URI-query strings.
\end{specnote}

\section{Management Interface}\label{sec:management-interface}

The full \omgref{AMQP-Management}{1.0} protocol is out-of-scope
for this specification (\S\,\ref{sec:out-of-scope}). However, an
endpoint without any management surface is opaque to standard
AMQP operations tooling. This section defines a small,
read-only management surface that an Endpoint Profile
implementation \rfckw{SHALL} expose.

\subsection{Catalog Address}

The address-catalog (\S\,\ref{sec:pim-discovery-bridging}) at
the source-address \texttt{\$catalog} is the primary management
read-channel. Operators inspect the topic-to-DDS-mapping by
attaching a Receiver link to this address. The catalog is
spec-mandatory.

\subsection{Metrics Address}\label{sec:management-metrics}

An endpoint \rfckw{SHALL} expose runtime metrics at the
source-address \texttt{\$metrics}. A Receiver link attached to
\texttt{\$metrics} \rfckw{SHALL} produce a stream of AMQP
messages, one per metric, where the body is an AMQP
\texttt{map} with the following keys:

\begin{itemize}
  \item \texttt{name} (string): metric identifier.
  \item \texttt{value} (long or double): current value.
  \item \texttt{unit} (symbol): unit-of-measure
        (e.g.\ \texttt{"count"}, \texttt{"bytes"},
        \texttt{"milliseconds"}, \texttt{"per-second"}).
  \item \texttt{timestamp} (timestamp): observation time.
\end{itemize}

\subsubsection{Mandatory Metric Names}

To permit standardised operations tooling, an endpoint
\rfckw{SHALL} emit at least the following metrics:

\begin{longtable}{p{4.6cm}p{2cm}p{5.5cm}}
  \toprule
  \textbf{name} & \textbf{unit} & \textbf{Meaning} \\
  \midrule
  \endhead
  \texttt{connections.active}        & \texttt{count}      & current open AMQP connections \\
  \texttt{connections.total}         & \texttt{count}      & cumulative connections accepted since start \\
  \texttt{transfers.received}        & \texttt{count}      & cumulative transfers received from peers \\
  \texttt{transfers.sent}            & \texttt{count}      & cumulative transfers sent to peers \\
  \texttt{transfers.unsettled}       & \texttt{count}      & current unsettled deliveries on all links \\
  \texttt{transfers.rate}            & \texttt{per-second} & sliding-window throughput \\
  \texttt{errors.decode}             & \texttt{count}      & cumulative \texttt{amqp:decode-error} occurrences \\
  \texttt{errors.unauthorized}       & \texttt{count}      & cumulative \texttt{amqp:unauthorized-access} occurrences \\
  \texttt{topics.exposed}            & \texttt{count}      & current address-catalog entries \\
  \texttt{transfers.dropped.loop}    & \texttt{count}      & cumulative samples dropped by bridge-loop detection (\S\,\ref{sec:bridge-coexistence}) \\
  \texttt{transfers.dropped.hop-cap} & \texttt{count}      & cumulative samples dropped because \texttt{dds:bridge-hop} exceeded the implementation cap \\
  \texttt{transfers.dropped.malformed-reply} & \texttt{count} & cumulative RPC replies rejected by Annex~D \S\,D.4 reply-validation \\
  \texttt{rpc.calls.timed-out}       & \texttt{count}      & cumulative RPC-aware calls that exceeded \texttt{rpc\_timeout\_ms} (\S\,D.4) \\
  \texttt{transfers.dropped.reconnect-overflow} & \texttt{count} & cumulative samples evicted by \texttt{KEEP\_LAST} HISTORY during a Bridge-Profile reconnect window (\S\,\ref{sec:reconnect}) \\
  \bottomrule
\end{longtable}

Implementations \rfckw{MAY} additionally emit vendor-specific
metrics with names prefixed \texttt{vendor:}; non-prefixed
metric names are reserved for this specification.

\subsection{Limitations}

\texttt{\$catalog} and \texttt{\$metrics} are the only
\emph{operational} management surfaces specified by this
document. The \texttt{\$audit} channel
(\S\,\ref{sec:audit-channel}) is a separate \emph{security}
surface, governed by Chapter~\ref{ch:security}; it is not part
of the operational management surface defined here.
Implementations \rfckw{MAY} expose additional vendor-prefixed
addresses (\texttt{\$vendor:\ldots}) for write-side management,
but such extensions are not portable. Deployments requiring
full read/write management \rfckw{SHALL} layer
\omgref{AMQP-Management}{1.0} on top of this specification as a
separate stack.

\begin{specnote}
\omgref{AMQP-Management}{1.0} is not the de-facto operations
interface of the AMQP-1.0-native brokers this specification
targets: RabbitMQ, Solace PubSub+, Microsoft Azure Service Bus,
Apache Qpid Broker-J, and Apache ActiveMQ Artemis each expose
broker-grade observability via vendor-specific HTTP/REST/JMX
endpoints rather than via AMQP-Management-1.0. Building a full
AMQP-Management-1.0 facade onto a DDS-AMQP endpoint would not
increase interoperability with industry-standard ops tooling;
broker-grade observability is delegated to the colocated broker
(Bridge Profile) or to broker-native management plus the
\texttt{\$metrics} stream defined in
\S\,\ref{sec:management-metrics}. This is a deliberate
scope-decision, not a deferred feature.
\end{specnote}

\section{Implementation Limits}\label{sec:pim-impl-limits}

To prevent denial-of-service via pathological inputs, an
endpoint implementation \rfckw{SHALL} enforce the following
upper bounds:

\begin{itemize}
  \item \textbf{Compound nesting depth}: the number of nested
        list/map/array values within a single AMQP body
        \rfckw{SHALL NOT} exceed an implementation-defined limit;
        the limit \rfckw{SHALL} be at least 16 and \rfckw{SHOULD}
        be at most 64.
  \item \textbf{Frame size}: the encoded size of a single AMQP
        frame \rfckw{SHALL NOT} exceed the connection's negotiated
        \texttt{max-frame-size}, which itself \rfckw{MUST} be at
        least 512 octets per
        \omgref{amqp-1.0-transport}{\S\,2.7.1.1}.
  \item \textbf{Catalog entries}: the number of entries in
        the address catalog (\S\,\ref{sec:pim-discovery-bridging})
        \rfckw{SHALL NOT} exceed an implementation-defined upper
        bound; the bound \rfckw{SHOULD} be at least 1024.
  \item \textbf{Unsettled deliveries}: the number of unsettled
        deliveries pending on a single receiver-link
        \rfckw{SHALL NOT} exceed
        $\min(\texttt{HISTORY.depth}, 65536)$ for KEEP\_LAST QoS
        and \rfckw{SHALL NOT} exceed an implementation-defined
        bound for KEEP\_ALL.
\end{itemize}

When any of the above limits is exceeded, the endpoint
\rfckw{SHALL} emit a structured error: for body-decoding limits
(\texttt{nesting depth}), a transfer \rfckw{SHALL} be released
with \texttt{rejected} disposition carrying error condition
\texttt{amqp:decode-error}. For session-level limits (catalog
size, unsettled deliveries), the connection \rfckw{SHALL} be
closed with \texttt{amqp:resource-limit-exceeded}.

% ============================================================================
%  Chapter 8 — Wire Encoding (PSM).
% ============================================================================
\chapter{Wire Encoding (Normative)}\label{ch:psm}

This chapter specifies how the platform-independent mappings of
Chapter~\ref{ch:pim} are realised on the AMQP wire.

\section{Body Encoding Modes}

Three body-encoding modes are defined. An Endpoint or Bridge
implementation \rfckw{MUST} support at least one mode and
\rfckw{SHOULD} support all three. The active mode per topic
is determined by the configuration mechanism
of~\S\,\ref{ch:configuration}.

\subsection{Pass-Through Octet-String Mode}\label{sec:psm-passthrough}

In Pass-Through mode, the body of an AMQP message
\rfckw{SHALL} be a single \texttt{data} section
(\omgref{amqp-1.0-messaging}{\S\,3.2.6}) whose octet content is the
serialised XCDR2 representation of the DDS sample.

The Properties section's \texttt{content-type} field
\rfckw{SHALL} be set to
\texttt{application/vnd.dds.xcdr2}. The Properties section's
\texttt{content-encoding} field \rfckw{MAY} be omitted.

\begin{specnote}
Pass-Through is the simplest and most efficient mode: the
endpoint performs no body translation. It \rfckw{REQUIRES} that
both AMQP peers understand XCDR2 and is therefore most useful in
DDS-AMQP-to-DDS-AMQP federation scenarios.
\end{specnote}

\subsection{JSON Mapping Mode}\label{sec:psm-json}

In JSON mode, the body of an AMQP message \rfckw{SHALL} be a single
\texttt{data} section whose octet content is a UTF-8-encoded JSON
document (RFC 8259) representing the DDS sample.

The Properties section's \texttt{content-type} field
\rfckw{SHALL} be set to \texttt{application/json}.

The mapping from IDL to JSON \rfckw{SHALL} follow these rules:

\begin{itemize}
  \item Primitive types: as defined by RFC 8259.
        Floating-point \texttt{NaN} and infinities are not
        representable in JSON and \rfckw{SHALL} cause the encoder
        to raise an error.
  \item \texttt{string} and \texttt{wstring}: JSON string.
  \item \texttt{octet} sequences: Base64 (RFC 4648) string.
  \item \texttt{enum}: JSON string of the symbolic name.
  \item \texttt{struct}: JSON object whose keys are the IDL
        field names. Optional fields with absent values
        \rfckw{SHOULD} be omitted (rather than encoded as
        \texttt{null}).
  \item \texttt{union}: JSON object with two keys:
        \texttt{discriminator} (the value) and \texttt{value}
        (the active branch). The \texttt{value} key
        \rfckw{SHALL} be omitted for empty branches.
  \item \texttt{sequence} and array: JSON array.
  \item \texttt{map}: JSON object.
  \item \texttt{Time\_t}: JSON object \texttt{\{"sec": $n$,
        "nanosec": $m$\}}.
  \item 16-byte identifiers (\texttt{InstanceHandle\_t},
        \texttt{GUID\_t}, X-Types \texttt{TypeIdentifier} when the
        identifier is 16 bytes): JSON string of 32 lowercase
        hexadecimal characters. \texttt{GUID\_t} \rfckw{MAY}
        alternatively be expressed in canonical RFC-4122 form
        (\texttt{8-4-4-4-12}) when the GUID is known to be
        RFC-4122-conformant. The \texttt{message-id} per-sample
        identifier (24 bytes, \S\,\ref{sec:psm-properties}) is
        encoded as a JSON string of 48 lowercase hexadecimal
        characters when surfaced through JSON mode.
\end{itemize}

\subsection{AMQP-Native Typed Mapping Mode}\label{sec:psm-amqp-native}

In AMQP-Native mode, the body of an AMQP message \rfckw{SHALL} be a
single \texttt{amqp-value} section
(\omgref{amqp-1.0-messaging}{\S\,3.2.8}) containing the DDS
sample encoded as the described composite of
\S\,\ref{sec:pim-composite-mapping}.

The Properties section's \texttt{content-type} field
\rfckw{MAY} be omitted. When present it
\rfckw{SHALL} be \texttt{application/vnd.dds.amqp-native}.

\section{Properties-Section Mapping}\label{sec:psm-properties}

The Properties section
(\omgref{amqp-1.0-messaging}{\S\,3.2.4}) of every transferred
message \rfckw{SHALL} be populated as follows:

\begin{longtable}{p{4cm}p{8cm}}
  \toprule
  \textbf{Field} & \textbf{DDS Source / Constraint} \\
  \midrule
  \endhead
  \texttt{message-id}     & per-sample unique identifier, encoded as \texttt{binary} (24 bytes), formed by concatenating the writer's RTPS GUID (16 bytes) and the sample's RTPS sequence-number (8 bytes, big-endian). The DDS \texttt{InstanceHandle\_t} is \emph{not} carried in \texttt{message-id} (it is a per-instance identifier and would mis-trigger broker-side duplicate-detection on consecutive samples of the same key); it is carried in the \texttt{dds:instance-handle} application-property (\S\,\ref{sec:psm-app-properties}). When \texttt{rpc\_aware = true} the per-sample identifier is overridden by the RPC-aware mapping defined in Annex~D \S\,D.1. \\
  \texttt{user-id}        & \rfckw{MAY} be set to the DDS-Security identity certificate fingerprint. \\
  \texttt{to}             & target address of the link. \\
  \texttt{subject}        & \rfckw{MAY} carry an application-defined routing-key or sub-filter. The DDS topic name is already conveyed by the link's source/target address (\S\,\ref{sec:pim-address-resolution}); \texttt{subject} \rfckw{MUST NOT} be required to duplicate it and \rfckw{SHALL} remain available to applications. \\
  \texttt{reply-to}       & \rfckw{MAY} be set for request-reply patterns. \\
  \texttt{correlation-id} & \rfckw{MAY} be used for application-level correlation; not specified by this document. \\
  \texttt{content-type}   & MIME type of the body, see~\S\,\ref{sec:psm-passthrough}--\ref{sec:psm-amqp-native}. \\
  \texttt{content-encoding} & \rfckw{SHOULD} be omitted; bodies are not transport-encoded. \\
  \texttt{absolute-expiry-time} & DDS LIFESPAN-derived expiry, when applicable. \\
  \texttt{creation-time}  & DDS source-timestamp of the sample. \\
  \texttt{group-id}       & \rfckw{SHALL} carry the canonical key string of the DDS instance per \S\,\ref{sec:pim-key-instance}, when the topic type has \texttt{@key} fields. \rfckw{SHALL} be omitted for unkeyed topics. \\
  \texttt{group-sequence} & \rfckw{MAY} carry an application-defined per-instance sequence-number; not constrained by this document. \\
  \texttt{reply-to-group-id} & \rfckw{MAY} be used by application; not specified by this document. \\
  \bottomrule
\end{longtable}

\section{Application-Properties Mapping}\label{sec:psm-app-properties}

The Application-Properties section
(\omgref{amqp-1.0-messaging}{\S\,3.2.5}) \rfckw{MAY} be used to
carry sample metadata that is not representable in the standard
Properties section. The following keys are defined by this
specification and \rfckw{SHALL} be used when applicable:

\begin{longtable}{p{4cm}p{8cm}}
  \toprule
  \textbf{Key} & \textbf{Meaning} \\
  \midrule
  \endhead
  \texttt{dds:nsec}       & nanoseconds component of \texttt{Time\_t} when \texttt{creation-time} preserves only milliseconds. \\
  \texttt{dds:partition}  & DDS Partition QoS value: an AMQP \texttt{list} of \texttt{string} (one element per partition) or, for the single-valued case, a single \texttt{string}. Conflict resolution against URI-query-form is governed by \S\,\ref{sec:pim-conflict-resolution}. \\
  \texttt{dds:domain-id}  & DDS Domain ID as integer. \\
  \texttt{dds:type-id}    & XTypes TypeIdentifier hex string (full 14-byte form). \\
  \texttt{dds:source-guid} & DDS-Security GUID\_t of the originating endpoint, hex-encoded. \\
  \texttt{dds:instance-handle} & DDS \texttt{InstanceHandle\_t} of the sample's instance, encoded as \texttt{binary} (16 bytes). Identifies the keyed instance, not the individual sample (which is identified by \texttt{message-id}). \\
  \texttt{dds:lifespan-ms} & remaining LIFESPAN duration in milliseconds. \\
  \texttt{dds:sample-state} & one of \texttt{read} / \texttt{not-read}. \\
  \texttt{dds:view-state}   & one of \texttt{new} / \texttt{not-new}. \\
  \texttt{dds:instance-state} & one of \texttt{alive} / \texttt{not-alive-disposed} / \texttt{not-alive-no-writers}. \\
  \texttt{dds:operation}  & one of \texttt{write} / \texttt{register} / \texttt{unregister} / \texttt{dispose}, signalling DDS instance-lifecycle action; see \S\,\ref{sec:pim-instance-lifecycle}. Default \texttt{write} when the property is absent. \\
  \texttt{dds:bridge-id}  & comma-separated list of bridge identifiers (RFC-4122 UUIDs) that have already forwarded this sample. See \S\,\ref{sec:bridge-coexistence}. \\
  \texttt{dds:bridge-hop} & non-negative integer, the count of bridge forwards this sample has undergone. See \S\,\ref{sec:bridge-coexistence}. \\
  \bottomrule
\end{longtable}

Application-defined keys \rfckw{SHALL NOT} use the \texttt{dds:}
prefix; this prefix is reserved for keys defined by future
revisions of this specification.

\section{Bridge Coexistence and Loop Prevention}\label{sec:bridge-coexistence}

A DDS-AMQP implementation \rfckw{MAY} coexist on the same DDS
domain with other forwarding agents (further DDS-AMQP
implementations, MQTT-DDS bridges, DDS-Web gateways,
vendor-specific bridges). Without explicit loop-prevention, a
sample re-published by a peer agent can re-enter the originating
implementation through DDS and be forwarded a second time,
producing exponential message duplication.

Every Endpoint Profile and Bridge Profile implementation
\rfckw{SHALL} apply the following rules to every sample it
forwards from DDS to AMQP and from AMQP to DDS. The Codec and
Codec-Lite Profiles do not implement forwarding and are not
bound by this section.

\begin{enumerate}
  \item \textbf{Bridge-id assignment.} The implementation
        \rfckw{SHALL} possess a single stable identifier of type
        \texttt{uuid} (RFC 4122), formed at startup and retained
        for the lifetime of the process — \emph{not} per topic.
        The identifier \rfckw{SHALL} be configured in the
        \texttt{bridge\_id} field of \texttt{AmqpEndpointConfig}
        or \texttt{AmqpBridgeConfig} (Annex~A) or generated as a
        type-4 UUID if absent. A process \rfckw{SHALL NOT} use
        different \texttt{bridge\_id} values for different
        topics; doing so defeats the drop-on-self-tag rule when a
        sample published via topic A re-enters via topic B.
  \item \textbf{Outbound stamp.} When forwarding a sample to AMQP,
        the implementation \rfckw{SHALL} append its
        \texttt{bridge\_id} to the \texttt{dds:bridge-id}
        application-property and \rfckw{SHALL} increment
        \texttt{dds:bridge-hop} by one (creating the property as
        \texttt{1} if absent).
  \item \textbf{Inbound drop.} When receiving a sample from AMQP
        for forwarding to DDS, the implementation \rfckw{SHALL}
        inspect \texttt{dds:bridge-id}; if its own
        \texttt{bridge\_id} appears in the list, the sample
        \rfckw{SHALL} be dropped silently (no DDS publication, no
        error to peer) and \texttt{transfers.dropped.loop}
        (\S\,\ref{sec:management-metrics}) \rfckw{SHALL} be
        incremented.
  \item \textbf{Hop-count cap.} The implementation \rfckw{SHALL}
        enforce an upper bound on \texttt{dds:bridge-hop}; the
        bound \rfckw{SHALL} be at least 8 and \rfckw{SHOULD} be
        at most 16 unless an operator explicitly raises it.
        Samples exceeding the cap \rfckw{SHALL} be dropped and
        \texttt{transfers.dropped.hop-cap} \rfckw{SHALL} be
        incremented.
\end{enumerate}

\begin{specnote}
DDS-side coordination (topic naming, \texttt{PARTITION} QoS,
\texttt{GROUP\_DATA}/\texttt{TOPIC\_DATA}) remains a useful
defence-in-depth supplement but is not the primary
loop-prevention mechanism; the rules above ensure correctness
even when DDS-side coordination is misconfigured.
Cross-protocol bridges that wish to participate in the same
loop-prevention scheme \rfckw{SHOULD} adopt the
\texttt{dds:bridge-id}/\texttt{dds:bridge-hop} convention.
\end{specnote}

% ============================================================================
%  Chapter 9 — Configuration.
% ============================================================================
\chapter{Configuration (Normative)}\label{ch:configuration}

\section{IDL-defined Mapping Schema}

The configuration of a DDS-AMQP endpoint or bridge \rfckw{SHALL}
conform to the IDL schema defined in Annex~A. Implementations
\rfckw{MAY} accept configuration in any concrete representation
(TOML, JSON, XML, native runtime API) provided the resulting
configuration is semantically equivalent to a value of the IDL
schema.

\section{XML Configuration}

Implementations \rfckw{MAY} accept configuration in the XML form
defined by \omgref{OMG DDS-XML}{\S\,5}, extended with the
\texttt{<dds-amqp>} element whose content matches the IDL schema
of Annex~A. The XML namespace
\texttt{http://www.zerodds.org/dds-amqp/v1.0}
\rfckw{SHALL} be associated with the \texttt{<dds-amqp>} element.

% ============================================================================
%  Chapter 10 — Security.
% ============================================================================
\chapter{Security (Normative)}\label{ch:security}

\section{TLS Bracket}\label{sec:security-tls}

When TLS is used, the connection \rfckw{SHALL} conform to
RFC~8446 (TLS 1.3). The endpoint \rfckw{SHALL} accept TLS 1.3
client connections; it \rfckw{MAY} additionally accept TLS 1.2
clients for backwards compatibility. The endpoint \rfckw{SHALL NOT}
accept TLS versions older than 1.2.

The endpoint \rfckw{SHALL} support the cipher suites mandated by
RFC~8446, \S\,9.1: \texttt{TLS\_AES\_128\_GCM\_SHA256},
\texttt{TLS\_AES\_256\_GCM\_SHA384}, and
\texttt{TLS\_CHACHA20\_POLY1305\_SHA256}.

The endpoint \rfckw{SHOULD} support TLS Server Name Indication (SNI,
RFC~6066, \S\,3) to permit virtual hosting of multiple endpoints
behind a single TCP listener.

Mutual TLS authentication (mTLS) \rfckw{MAY} be required by
configuration. When required, the endpoint \rfckw{SHALL} reject
connections whose client certificate is absent, expired, or not
issued by a trusted CA.

\section{SASL Mechanisms}\label{sec:security-sasl}

Endpoint and Bridge Profile implementations \rfckw{SHALL} support
the following SASL mechanisms (RFC 4422):

\begin{itemize}
  \item \texttt{PLAIN} (RFC 4616): username and password
        authentication.
  \item \texttt{ANONYMOUS} (RFC 4505): unauthenticated access. The
        configuration controls whether \texttt{ANONYMOUS} is
        offered.
  \item \texttt{EXTERNAL} (RFC 4422 \S\,7): authentication
        delegated to the underlying transport, typically via mTLS
        client certificate.
\end{itemize}

Implementations \rfckw{MAY} additionally support:

\begin{itemize}
  \item \texttt{SCRAM-SHA-256} (RFC 5802): salted challenge-response
        authentication.
  \item Vendor-specific mechanisms in conformance with RFC 4422
        \S\,3.
\end{itemize}

\subsection{Mandatory TLS for SASL PLAIN}\label{sec:security-sasl-plain-tls}

SASL PLAIN transmits the user-name and password as cleartext on
the wire (RFC 4616 \S\,2). Implementations \rfckw{MUST NOT} offer
SASL PLAIN as an Endpoint Profile mechanism when the underlying
transport is unencrypted, and \rfckw{MUST NOT} initiate SASL
PLAIN as a Bridge Profile mechanism unless the connection has
completed a successful TLS handshake first.

A conformant endpoint \rfckw{SHALL} reject any client SASL-init
frame carrying mechanism \texttt{PLAIN} on an unencrypted
connection with the SASL outcome \texttt{auth} (RFC 4422
\S\,3.5), and \rfckw{SHALL} immediately close the connection.

A conformant bridge \rfckw{SHALL NOT} send a SASL-init frame
with mechanism \texttt{PLAIN} unless its TCP connection is
wrapped in TLS as defined in \S\,\ref{sec:security-tls}. The
bridge \rfckw{MAY} negotiate \texttt{ANONYMOUS} or
\texttt{EXTERNAL} on unencrypted transports.

\section{Cross-Reference to DDS-Security}\label{sec:security-dds-security}

This specification does not modify
\omgref{OMG DDS-Security}{\S\,1.4}. A DDS-AMQP endpoint
\rfckw{MAY} simultaneously employ DDS-Security plugins on the DDS
side and AMQP-layer security (TLS+SASL) on the AMQP side; the two
are independent in their cryptographic operation but coupled in
their identity model as defined below.

\subsection{Outbound Signing and Encryption}

When DDS-Security is active and an AMQP-derived sample is forwarded
into the DDS domain, the endpoint \rfckw{SHALL} sign and encrypt
the outbound DDS submessage according to the configured DDS
plugins. This step is governed entirely by
\omgref{OMG DDS-Security}{\S\,7-\S\,9} and not modified here.

\subsection{Identity Mapping (normative)}\label{sec:security-identity-mapping}

For multi-vendor interoperability the AMQP-side authentication
identity \rfckw{SHALL} be translated into a DDS-Security
IdentityToken using the deterministic mapping below. Earlier
drafts of this specification left this mapping
\textsc{implementation-defined}; that wording has been retracted
as of v1.0.0-beta1 because it caused federation failures between
bridges of different vendors.

\begin{longtable}{p{3.5cm}p{9cm}}
  \toprule
  \textbf{SASL Mechanism} & \textbf{IdentityToken construction} \\
  \midrule
  \endhead
  \texttt{PLAIN}           & \texttt{class\_id = "zerodds:Auth:SASL-Username:1.0"}; the AMQP \texttt{authcid} (username) \rfckw{SHALL} be inserted into the token's \texttt{subject\_name} property as the literal value \texttt{"CN=" + authcid}. The \texttt{certificate} property is absent by construction; permission evaluation \rfckw{SHALL} use an allow-list keyed by \texttt{subject\_name}. The \texttt{DDS:Auth:PKI-DH:1.0} class\_id is \emph{not} used: that OMG class\_id is bound to a PKI-DH-with-X.509 plugin contract which would reject username-only tokens. \\
  \texttt{ANONYMOUS}       & \texttt{class\_id = "zerodds:Auth:Anonymous:1.0"}; \texttt{subject\_name} \rfckw{SHALL} be set to the literal string \texttt{"CN=ANONYMOUS"}. Permission evaluation \rfckw{SHALL} treat anonymous identities as a single shared subject. \\
  \texttt{EXTERNAL}        & \texttt{class\_id = "DDS:Auth:PKI-DH:1.0"}; the X.509 client certificate from the underlying TLS handshake \rfckw{SHALL} be inserted verbatim into the \texttt{certificate} property. \texttt{subject\_name} \rfckw{SHALL} be set to the certificate's Subject DN as serialised by RFC~4514. This is the only mechanism that produces an OMG-PKI-DH-conformant IdentityToken. \\
  \texttt{SCRAM-SHA-256}   & \texttt{class\_id = "zerodds:Auth:SASL-SCRAM-SHA256:1.0"}; the post-authentication \texttt{authcid} \rfckw{SHALL} be inserted into \texttt{subject\_name} as \texttt{"CN=" + authcid}; \texttt{certificate} is absent. Permission evaluation \rfckw{SHALL} use an allow-list keyed by \texttt{subject\_name}. \\
  \bottomrule
\end{longtable}

\subsection{Permission Evaluation}

The constructed IdentityToken \rfckw{SHALL} be passed to the
configured DDS-Security AccessControl plugin's
\texttt{check\_create\_datawriter} /
\texttt{check\_create\_datareader} entry points
(\omgref{OMG DDS-Security}{\S\,8.4.2}). If the plugin returns
\texttt{NOT\_ALLOWED}, the endpoint \rfckw{SHALL} reject the
AMQP link attachment with link error
\texttt{amqp:unauthorized-access}.

\subsection{Determinism Across Vendors}

Two conformant DDS-AMQP endpoints in a multi-vendor federation
deployment (\S\,\ref{sec:topology-hybrid}, Use-Case
\S\,6.2.3) \rfckw{SHALL} produce byte-identical IdentityToken
\texttt{subject\_name} values for the same AMQP-side
authentication identity. This determinism is what permits a
single DDS-Security Permissions Document to govern both
endpoints.

\subsection{No-Bypass Guarantee}\label{sec:security-no-bypass}

When the DDS domain into which the endpoint forwards traffic has
DDS-Security activated, the endpoint or bridge \rfckw{SHALL NOT}
permit the AMQP client to publish or subscribe to that domain
without first having satisfied the DDS-Security permission
evaluation chain. Specifically:

\begin{enumerate}
  \item The AMQP-side authentication \rfckw{MUST} have completed
        successfully (SASL outcome \texttt{ok}).
  \item The constructed IdentityToken
        (\S\,\ref{sec:security-identity-mapping}) \rfckw{MUST}
        be presented to the configured DDS-Security
        AccessControl plugin's
        \texttt{check\_create\_datawriter} or
        \texttt{check\_create\_datareader} entry point.
  \item Only on a \texttt{ALLOWED} return value \rfckw{MAY} the
        endpoint create the corresponding DDS DataWriter or
        DataReader and forward AMQP traffic to it.
\end{enumerate}

A return value of \texttt{NOT\_ALLOWED} \rfckw{SHALL} cause the
endpoint to reject the AMQP link attachment with link error
\texttt{amqp:unauthorized-access}; the AMQP connection
\rfckw{MAY} remain open for further attachments to other topics
that the same identity is authorised for.

\subsubsection{Fallback when DDS-Security is Not Active}

When the DDS domain has DDS-Security \emph{not} configured (no
authentication, access-control, or cryptography plugin instances),
the AMQP-side authentication still constitutes the only access
control between the AMQP peer and the DDS domain. In this
configuration, the endpoint \rfckw{SHALL}:

\begin{enumerate}
  \item Treat a successful SASL outcome (\texttt{ok}) as
        sufficient authorisation for any topic mapping in its
        configured catalog.
  \item Reject any link attachment to an address not in the
        catalog with \texttt{amqp:not-found}, regardless of
        \texttt{permit\_dynamic\_topics} setting (DDS-Security-
        absent deployments \rfckw{MUST NOT} permit dynamic topic
        creation).
  \item Emit an audit record (per \S\,\ref{sec:audit-channel})
        for every successful link attachment, since no
        DDS-Security audit trail exists.
\end{enumerate}

This fallback preserves the No-Bypass property by making AMQP-
side authentication the binding security boundary; it does
\emph{not} create an unauthenticated path into the DDS domain.

\subsubsection{Audit Channel}\label{sec:audit-channel}

An endpoint \rfckw{SHALL} expose an audit-record stream at the
source-address \texttt{\$audit}. A Receiver link attached to
\texttt{\$audit} \rfckw{SHALL} produce a stream of AMQP messages,
one per audit event, where the body is an AMQP \texttt{map}
with the following keys:

\begin{itemize}
  \item \texttt{event} (symbol): event class —
        \texttt{"link.attach.success"},
        \texttt{"link.attach.rejected"},
        \texttt{"connection.opened"},
        \texttt{"connection.closed"},
        \texttt{"sasl.auth.failed"}.
  \item \texttt{timestamp} (timestamp): event time.
  \item \texttt{subject} (string): authenticated subject identifier
        (e.g.\ SASL authcid, X.509 subject DN, or
        \texttt{"ANONYMOUS"}).
  \item \texttt{address} (string): AMQP address of the link, when
        applicable.
  \item \texttt{detail} (string): free-text human-readable detail,
        bounded to 256 octets.
\end{itemize}

The audit stream is read-only; subscribers \rfckw{MUST} attach
Receiver links only. The endpoint \rfckw{MAY} additionally write
the same records to a local log; the AMQP-exposed stream is
the canonical machine-readable form for operations tooling.

\subsection{Governance Document Mapping}\label{sec:security-governance}

The DDS-Security Governance Document
(\omgref{OMG DDS-Security}{\S\,9.4.1.1}) governs which DDS
topics require which protections (discovery-encryption,
metaprotection, dataprotection, sign/encrypt). The endpoint
\rfckw{SHALL} apply the configured Governance Document
unchanged on the DDS side.

The AMQP side has its own protection layer (TLS + SASL); the
two are composable but not equivalent. The detailed rules for
mapping per-topic Governance protection-kinds to AMQP transport
requirements are normatively specified in
\S\,\ref{sec:per-link-governance}.

\subsection{No Implicit Trust Between Profiles}

Endpoint and Bridge profiles \rfckw{MUST NOT} treat the AMQP-side
authentication as a free pass to the DDS side. The
identity-translation chain
(\S\,\ref{sec:security-identity-mapping}) is the only authorised
bridge; implementations \rfckw{SHALL NOT} provide
``trusted-bridge'' modes that skip DDS-Security evaluation.

\subsection{Bridge-Profile Dual Identity}\label{sec:bridge-dual-identity}

A Bridge Profile implementation operates at the intersection of
two distinct authentication domains:

\begin{enumerate}
  \item the \emph{broker-side identity} — the bridge authenticates
        \emph{outbound} to an external broker (e.g.\ RabbitMQ)
        via SASL on its outgoing AMQP connection;
  \item the \emph{DDS-side identity} — the bridge participates as
        a DDS DomainParticipant inside the configured DDS
        domain, with its own DDS-Security IdentityToken
        (\omgref{OMG DDS-Security}{\S\,9.4.2.4}).
\end{enumerate}

The two identities \rfckw{SHALL NOT} be conflated. The bridge's
DDS-Security IdentityToken \rfckw{SHALL} be configured
independently from the broker-side credentials. Specifically:

\begin{itemize}
  \item The DDS-Security AccessControl plugin \rfckw{SHALL}
        evaluate the bridge's DDS IdentityToken (not the broker
        credentials) when deciding whether the bridge may
        publish/subscribe on a DDS topic.
  \item AMQP messages forwarded by the bridge into DDS
        \rfckw{SHALL NOT} carry the broker's IdentityToken; if
        per-message originator identity is required, the
        AMQP-side authcid \rfckw{SHALL} be carried in the
        Application-Property \texttt{dds:source-authcid}.
\end{itemize}

\subsection{Per-Link Governance Resolution}\label{sec:per-link-governance}

When the DDS-Security Governance Document
(\omgref{OMG DDS-Security}{\S\,9.4.1.1}) assigns different
protection-kinds to different topics within the same domain
(common in mixed-sensitivity deployments), the endpoint
\rfckw{SHALL} resolve TLS-requirement on a \emph{per-link}
basis, not per-connection:

\begin{enumerate}
  \item For every link attachment, the endpoint \rfckw{SHALL}
        consult the Governance entry for the target topic.
  \item If the Governance \texttt{rtps\_protection\_kind} or
        \texttt{data\_protection\_kind} is
        \texttt{ENCRYPT}-grade, the endpoint \rfckw{SHALL}
        require that the underlying AMQP connection has
        completed a TLS handshake; if not, the link attachment
        \rfckw{SHALL} be rejected with
        \texttt{amqp:unauthorized-access}.
  \item Otherwise (\texttt{NONE} or \texttt{SIGN}-only), the
        link \rfckw{MAY} be accepted on either encrypted or
        plain transport.
\end{enumerate}

A single AMQP connection \rfckw{MAY} therefore host both
ENCRYPT-grade and SIGN-only links provided the connection
itself satisfies the strictest requirement (TLS for ENCRYPT-grade).

\subsection{Reconnect Behaviour}\label{sec:reconnect}

After an \texttt{amqp:connection:framing-error},
\texttt{amqp:resource-limit-exceeded}, or any abrupt TCP-reset:

\begin{itemize}
  \item \textbf{Endpoint Profile} (mandatory): the endpoint
        \rfckw{SHALL} continue to accept new client connections;
        the failed connection's in-flight unsettled deliveries
        \rfckw{SHALL} be released per
        \omgref{amqp-1.0-transport}{\S\,3.4.6}.
  \item \textbf{Bridge Profile} (mandatory): the bridge
        \rfckw{SHALL} attempt reconnection to the configured
        upstream broker after any abnormal disconnect, and
        \rfckw{SHALL} re-attach all configured links and resume
        sample flow once the connection is re-established. While
        the reconnect attempt is in progress, the bridge
        \rfckw{SHALL} preserve DDS-side samples to the extent
        permitted by the link's configured HISTORY QoS:
        \texttt{KEEP\_ALL} retains all samples; \texttt{KEEP\_LAST(N)}
        retains the most recent $N$ per instance and evicts older
        samples on overflow as part of normal HISTORY policy.
        Samples evicted by \texttt{KEEP\_LAST} overflow during a
        disconnect window \rfckw{SHALL} be counted in the metric
        \texttt{transfers.dropped.reconnect-overflow}
        (\S\,\ref{sec:management-metrics}). The bridge
        \rfckw{SHALL NOT} drop samples for any reason other than
        the configured HISTORY policy or implementation
        resource-limits documented in \S\,\ref{sec:pim-impl-limits}.
  \item \textbf{Reconnect Pacing} (recommendation): the bridge
        \rfckw{SHOULD} apply exponential back-off to avoid
        connection-storms when the broker is unreachable.
        Reasonable defaults are: initial delay
        $\approx 1$\,s, multiplier $\approx 2$, cap at
        $\approx 60$\,s; deployments \rfckw{MAY} configure
        different values to match operational policy. The
        specification does not mandate exact timing values
        because pacing requirements are deployment-specific.
\end{itemize}

% ============================================================================
%  Chapter 11 — Error Conditions.
% ============================================================================
\chapter{Error Conditions (Normative)}\label{ch:error-conditions}

This chapter normatively maps DDS-AMQP failure modes to AMQP-1.0
error-condition symbols. Implementations \rfckw{SHALL} use the
exact symbols defined here so that diagnostic output and
automatic retry-logic on the peer side is interoperable.

\section{Encode/Decode Failures}

\begin{longtable}{p{4.0cm}p{8.0cm}}
  \toprule
  \textbf{Failure} & \textbf{AMQP error condition + scope} \\
  \midrule
  \endhead
  IDL type with no defined AMQP mapping (encoder-side) & \texttt{amqp:decode-error}; emitted as a \texttt{disposition} with \texttt{rejected} state on the offending transfer; the link \rfckw{SHALL} remain open. \\
  XCDR2 sample malformed (decoder-side) & \texttt{amqp:decode-error}; same scope as above. \\
  Compound nesting depth exceeds limit (\S\,\ref{sec:pim-impl-limits}) & \texttt{amqp:decode-error}; rejected disposition. \\
  String/wstring contains invalid UTF-8 / surrogate / out-of-range codepoint (\S\,\ref{sec:wstring-encoding}) & \texttt{amqp:decode-error}; rejected disposition. \\
  IDL char with no Unicode mapping in configured source code-page (\S\,\ref{sec:char-encoding-detail}) & \texttt{amqp:decode-error}; rejected disposition. \\
  Hash-truncation collision detected (descriptor matches but \texttt{dds:type-id} differs) & \texttt{amqp:decode-error}; rejected disposition. \\
  Frame size exceeds connection's \texttt{max-frame-size} & \texttt{amqp:connection:framing-error}; connection \rfckw{SHALL} be closed. \\
  \bottomrule
\end{longtable}

\section{Connection / Session / Link Lifecycle Failures}

\begin{longtable}{p{4.0cm}p{8.0cm}}
  \toprule
  \textbf{Failure} & \textbf{AMQP error condition + scope} \\
  \midrule
  \endhead
  Idle-timeout exceeded (\S\,\ref{sec:pim-impl-limits}) & \texttt{amqp:resource-limit-exceeded}; connection close. \\
  \texttt{max\_connections} cap reached on a new attach attempt & \texttt{amqp:resource-limit-exceeded}; connection close. \\
  Catalog-entries cap reached & \texttt{amqp:resource-limit-exceeded}; connection close. \\
  \texttt{permit\_dynamic\_topics = false} and address not in catalog (\S\,\ref{sec:pim-unknown-address}) & \texttt{amqp:not-found}; link error; connection remains open. \\
  Sender-link declares \texttt{TRANSIENT} or \texttt{PERSISTENT} durability (\S\,\ref{sec:pim-settlement}) & \texttt{amqp:not-implemented}; link error. \\
  \texttt{terminus.durable = unsettled-state} on inbound attach & \texttt{amqp:not-implemented}; link error. \\
  Unknown \texttt{dds:operation} value (\S\,\ref{sec:pim-instance-lifecycle}) & \texttt{amqp:not-implemented}; link error. \\
  DDS-Security AccessControl plugin returns \texttt{NOT\_ALLOWED} (\S\,\ref{sec:security-identity-mapping}) & \texttt{amqp:unauthorized-access}; link error. \\
  SASL PLAIN attempted on unencrypted transport (\S\,\ref{sec:security-sasl-plain-tls}) & SASL-layer outcome \texttt{auth} (RFC 4422); the AMQP layer is never reached, so no AMQP error-condition applies; the TCP connection is closed. \\
  Unknown described-composite descriptor & \texttt{amqp:not-implemented}; rejected disposition (per-transfer); link remains open. \\
  \bottomrule
\end{longtable}

\section{Instance-Lifecycle Failures}

\begin{longtable}{p{4.0cm}p{8.0cm}}
  \toprule
  \textbf{Failure} & \textbf{AMQP error condition + scope} \\
  \midrule
  \endhead
  \texttt{dds:operation = unregister} on unknown instance & \texttt{amqp:precondition-failed}; rejected disposition. \\
  \texttt{dds:operation = dispose} on unknown instance & \texttt{amqp:precondition-failed}; rejected disposition. \\
  \texttt{dds:operation = register} body lacks key fields & \texttt{amqp:decode-error}; rejected disposition. \\
  \bottomrule
\end{longtable}

\section{Diagnostic Description Strings}

Every error-condition emission \rfckw{SHOULD} carry a
human-readable description string in the \texttt{description}
field of the AMQP \texttt{error} composite
(\omgref{amqp-1.0-transport}{\S\,2.8.15}). Implementations
\rfckw{SHOULD NOT} place sensitive data (e.g.\ secret values,
authentication credentials) in the description.

% ============================================================================
%  Annex A — IDL Configuration Schema (normative).
% ============================================================================
\annex{A}{normative}{IDL Configuration Schema}

The following IDL fragment normatively defines the configuration
schema for DDS-AMQP implementations.

\begin{lstlisting}
module zerodds {
  module amqp {

    @final
    enum SaslMechanism {
      SASL_PLAIN,
      SASL_ANONYMOUS,
      SASL_EXTERNAL,
      SASL_SCRAM_SHA_256
    };

    @final
    enum BodyEncodingMode {
      MODE_PASSTHROUGH,
      MODE_JSON,
      MODE_AMQP_NATIVE
    };

    @final
    enum TimeMapping {
      MAPPING_STANDARD,    // AMQP timestamp + dds:nsec (default)
      MAPPING_COMPOSITE    // dds-time described composite (opt-in)
    };

    @final
    enum DescriptorForm {
      DESC_TRUNCATED,      // ulong, first 8 octets of equivalence-hash (default)
      DESC_FULL            // symbol, full 14-byte TypeIdentifier (collision-safe)
    };

    @final
    enum LinkDirection {
      DIR_PRODUCER_TO_DDS,
      DIR_DDS_TO_CONSUMER,
      DIR_BOTH
    };

    @final
    struct TopicMapping {
      string<256>      amqp_address;
      string<256>      dds_topic;
      string<256>      dds_type_name;
      uint32           dds_domain_id;          // default 0
      sequence<string<128>, 16> dds_partition;   // default empty sequence; each element is one PARTITION QoS string
      BodyEncodingMode mode;                   // default MODE_PASSTHROUGH
      TimeMapping      time_mapping;           // default MAPPING_STANDARD
      DescriptorForm   descriptor_form;        // default DESC_TRUNCATED
      boolean          rpc_aware;              // default false (see Annex D)
      uint32           rpc_timeout_ms;         // default 30000; only meaningful when rpc_aware=true
      LinkDirection    direction;              // default DIR_BOTH
    };

    @final
    struct TlsConfig {
      boolean enabled;
      string<512> cert_path;
      string<512> key_path;
      string<512> ca_path;
      boolean require_client_cert;
    };

    @final
    struct SaslConfig {
      sequence<SaslMechanism, 8> enabled_mechanisms;
      string<256> credential_store_uri;
    };

    @final
    struct ResourceLimits {
      uint32 max_connections;
      uint32 max_sessions_per_connection;
      uint32 max_links_per_session;
      uint32 max_frame_size;
      uint64 idle_timeout_ms;
    };

    @final
    struct DynamicTopicConfig {
      boolean        permit_dynamic_topics;
      string<256>    dynamic_topic_default_type;
      BodyEncodingMode default_mode;
    };

    @final
    struct AmqpEndpointConfig {
      @key string<128>       endpoint_name;
      string<256>            listen_uri;
      TlsConfig              tls;
      SaslConfig             sasl;
      sequence<TopicMapping> topics;
      DynamicTopicConfig     dynamic;
      ResourceLimits         limits;
      string<36>             bridge_id;             // RFC-4122 UUID; "" => generate type-4 at startup. Process-wide; SHALL be stable for the lifetime of the endpoint process. See sec:bridge-coexistence.
      uint8                  bridge_hop_cap;        // default 8, max 16 unless overridden. Process-wide.
    };

    @final
    struct AmqpBridgeConfig {
      @key string<128>       bridge_name;
      string<256>            upstream_uri;
      TlsConfig              tls;
      SaslConfig             sasl;
      sequence<TopicMapping> topics;
      string<36>             bridge_id;             // process-wide; same semantics as AmqpEndpointConfig.bridge_id
      uint8                  bridge_hop_cap;        // process-wide; same semantics as AmqpEndpointConfig.bridge_hop_cap
    };

  };  // module amqp
};  // module zerodds
\end{lstlisting}

% ============================================================================
%  Annex B — Examples (informative).
% ============================================================================
\annex{B}{informative}{Examples}

The configurations in this annex are written in TOML
(RFC-equivalent INI-style) as a \emph{concrete representation}
of the IDL schema of Annex~A. The mapping between the IDL
\texttt{struct} types of Annex~A and the TOML rendering is the
following:

\begin{itemize}
  \item Each top-level \texttt{struct} (e.g.\ \texttt{AmqpEndpointConfig})
        becomes a TOML table named after the lower-cased
        struct-name without the \texttt{Amqp} prefix
        (\texttt{[endpoint]}, \texttt{[bridge]}).
  \item Nested \texttt{struct}-typed fields become sub-tables
        (\texttt{[endpoint.tls]}, \texttt{[endpoint.sasl]},
        \texttt{[endpoint.dynamic]}, \texttt{[endpoint.limits]}).
  \item \texttt{sequence<TopicMapping>} becomes an array-of-tables
        (\texttt{[[endpoint.topics]]}).
  \item IDL \texttt{enum} values are rendered as TOML strings
        with the literal enumerator-name
        (e.g.\ \texttt{"MODE\_JSON"}, \texttt{"SASL\_PLAIN"}).
  \item IDL \texttt{boolean} renders as TOML \texttt{true}/\texttt{false};
        numeric types as TOML numbers; bounded strings verbatim.
\end{itemize}

This concrete representation is one of several permitted forms;
implementations \rfckw{MAY} accept JSON or XML
(\S\,\ref{ch:configuration}) as semantically equivalent
alternatives.

\section*{B.1 Edge Sensor to DDS Action Planner (Endpoint Profile)}

An edge sensor device emits sensor results via the
\texttt{pika} Python AMQP library. A ZeroDDS-AMQP endpoint listens
on \texttt{amqps://0.0.0.0:5671}, accepts the device's connection,
and writes the received samples to the DDS topic
\texttt{TrackingResult}.

\begin{lstlisting}
# TOML configuration (example concrete representation of the IDL
# schema in Annex A).
[endpoint]
endpoint_name = "edge-gateway-1"
listen_uri    = "amqps://0.0.0.0:5671"

[endpoint.tls]
enabled              = true
cert_path            = "/etc/zerodds/cert.pem"
key_path             = "/etc/zerodds/key.pem"
require_client_cert  = true

[endpoint.sasl]
enabled_mechanisms = ["SASL_PLAIN", "SASL_EXTERNAL"]

[[endpoint.topics]]
amqp_address  = "tracking"
dds_topic     = "TrackingResult"
dds_type_name = "robotics::TrackingResult"
dds_domain_id = 0
mode          = "MODE_JSON"
direction     = "DIR_PRODUCER_TO_DDS"

[endpoint.limits]
max_connections             = 1024
max_sessions_per_connection = 8
max_links_per_session       = 64
max_frame_size              = 65536
idle_timeout_ms             = 60000
\end{lstlisting}

\section*{B.2 DDS to Cloud Subscriber via External Broker (Bridge Profile)}

A ZeroDDS-AMQP bridge connects to RabbitMQ as both producer and
consumer, forwarding the DDS topic \texttt{Telemetry} to the AMQP
exchange \texttt{telemetry.fanout}.

\begin{lstlisting}
[bridge]
bridge_name  = "telemetry-bridge"
upstream_uri = "amqps://gw.zerodds.example:5671/telemetry"

[bridge.tls]
enabled  = true
ca_path  = "/etc/zerodds/ca-bundle.pem"

[bridge.sasl]
enabled_mechanisms   = ["SASL_PLAIN"]
credential_store_uri = "file:///etc/zerodds/secrets/rabbitmq.toml"

[[bridge.topics]]
amqp_address  = "telemetry.fanout"
dds_topic     = "Telemetry"
dds_type_name = "platform::TelemetrySample"
dds_domain_id = 7
mode          = "MODE_PASSTHROUGH"
direction     = "DIR_DDS_TO_CONSUMER"
\end{lstlisting}

\section*{B.3 Multi-Vendor Federation}

Two ZeroDDS-AMQP endpoints, each in a separate DDS domain, federate
via a directly-attached AMQP link. Endpoint~A exposes domain~3's
\texttt{Alerts} topic as the AMQP address \texttt{alerts}; endpoint~B
attaches a Receiver link to A's \texttt{alerts} address and writes
the samples into domain~5's \texttt{IngressAlerts} topic.

The catalog mechanism of~\S\,\ref{sec:pim-discovery-bridging} permits
endpoint~B to discover endpoint~A's catalog by attaching to
\texttt{\$catalog}, listing all available addresses, and selecting
\texttt{alerts} programmatically.

% ============================================================================
%  Annex C — Compliance Test Suite (normative).
% ============================================================================
\annex{C}{normative}{Compliance Test Suite}

This annex enumerates the test cases an implementation
\rfckw{MUST} satisfy to claim conformance to a given profile.

\section*{C.1 Endpoint-Profile Tests}

\begin{enumerate}
  \item \textbf{C.1.1 Connection Open (TLS+PLAIN)}: an AMQP-1.0
        client \rfckw{SHALL} be able to complete a TLS-1.3
        handshake to the endpoint, then authenticate via SASL
        PLAIN, exchange \texttt{open}/\texttt{open} performatives,
        and remain connected for at least one second.
  \item \textbf{C.1.2 SASL PLAIN Rejection on Plain Transport}:
        an attempt by an AMQP-1.0 client to authenticate via
        SASL PLAIN over an unencrypted (non-TLS) transport
        \rfckw{SHALL} be rejected with SASL outcome \texttt{auth}
        and the underlying TCP connection \rfckw{SHALL} be closed,
        per \S\,\ref{sec:security-sasl-plain-tls}.
  \item \textbf{C.1.3 Sender-Link Producer-to-DDS}: an
        AMQP-1.0 client \rfckw{SHALL} be able to attach a Sender
        link to a configured topic address and have its transferred
        sample appear on the corresponding DDS topic, observable by
        a DDS DataReader.
  \item \textbf{C.1.4 Receiver-Link DDS-to-Consumer}: an AMQP-1.0
        client \rfckw{SHALL} be able to attach a Receiver link to
        a configured topic address and receive subsequent
        DDS-published samples on the same topic.
  \item \textbf{C.1.5 Settlement-Mode Reliable}: a DDS DataWriter
        with \texttt{RELIABLE} reliability \rfckw{SHALL} produce
        unsettled AMQP transfers; the endpoint \rfckw{SHALL}
        only acknowledge completion to the DDS layer after
        receiving an \texttt{accepted} disposition.
  \item \textbf{C.1.6 Settlement-Mode Best-Effort}: a DDS
        DataWriter with \texttt{BEST\_EFFORT} reliability
        \rfckw{SHALL} produce settled (pre-settled) AMQP transfers.
  \item \textbf{C.1.7 Address-Resolution Wildcard}: an
        attached link with a wildcard address \rfckw{SHALL}
        receive transfers from all matching DDS Partitions.
  \item \textbf{C.1.8 Catalog-Address}: an AMQP-1.0 client
        \rfckw{SHALL} be able to attach a Receiver link to
        \texttt{\$catalog} and receive at least one entry per
        configured topic mapping.
  \item \textbf{C.1.9 Idle-Timeout}: a connection with no traffic
        for the configured idle-timeout duration \rfckw{SHALL} be
        closed by the endpoint with a \texttt{close} performative
        carrying error condition
        \texttt{amqp:resource-limit-exceeded} or
        \texttt{amqp:connection:framing-error}.
  \item \textbf{C.1.10 Concurrent Connections}: the endpoint
        \rfckw{SHALL} accept at least
        \texttt{max\_connections}/2 simultaneous connections and
        serve them without rejection.
  \item \textbf{C.1.11 Pass-Through Encoding}: a sample written by
        the endpoint with \texttt{MODE\_PASSTHROUGH} \rfckw{SHALL}
        round-trip back to a byte-identical XCDR2 representation
        when consumed by a DDS-aware AMQP client.
  \item \textbf{C.1.12 JSON Encoding}: a sample written by the
        endpoint with \texttt{MODE\_JSON} \rfckw{SHALL} be valid
        JSON (RFC 8259) and decodable by a JSON-aware AMQP
        client into the original IDL value.
  \item \textbf{C.1.13 Metrics Link}: an AMQP-1.0 client
        \rfckw{SHALL} be able to attach a Receiver link to
        \texttt{\$metrics}; within an observation window of
        $\le$\,30\,s the link \rfckw{SHALL} deliver at least one
        sample for every Mandatory Metric enumerated in
        \S\,\ref{sec:management-metrics} (the test \rfckw{MAY}
        induce traffic to ensure non-zero counters are emitted
        for the throughput-related metrics). Each sample
        \rfckw{SHALL} carry the AMQP-\texttt{map} body shape
        defined in \S\,\ref{sec:management-metrics}.
  \item \textbf{C.1.14 Audit Link}: an AMQP-1.0 client
        \rfckw{SHALL} be able to attach a Receiver link to
        \texttt{\$audit}; a second AMQP-1.0 client subsequently
        opens a connection and completes SASL successfully on
        a different link. The audit-Receiver \rfckw{SHALL}
        observe a \texttt{link.attach.success} (or equivalent
        per \S\,\ref{sec:audit-channel}) event whose audit
        record carries the second client's authenticated
        \texttt{subject\_name}.
  \item \textbf{C.1.15 Loop-Prevention (Self-Tag and Hop-Cap)}:
        two sub-tests, both \rfckw{SHALL} pass.
        (a)~The test attaches a Sender link and writes a
        sample whose \texttt{dds:bridge-id} application-property
        contains the endpoint's own configured \texttt{bridge\_id}
        (\S\,\ref{sec:bridge-coexistence}). The endpoint
        \rfckw{SHALL} drop the sample without DDS publication
        and \rfckw{SHALL} increment
        \texttt{transfers.dropped.loop} by one (observable via
        C.1.13).
        (b)~The test writes a sample whose \texttt{dds:bridge-hop}
        is one greater than the endpoint's configured
        \texttt{bridge\_hop\_cap}. The endpoint \rfckw{SHALL}
        drop the sample and \rfckw{SHALL} increment
        \texttt{transfers.dropped.hop-cap} by one.
\end{enumerate}

\section*{C.2 Bridge-Profile Tests}

\begin{enumerate}
  \item \textbf{C.2.1 Outbound Connection}: the bridge
        \rfckw{SHALL} establish a connection to a configured
        upstream broker.
  \item \textbf{C.2.2 Sender-Link Bridge-to-Broker}: the bridge
        \rfckw{SHALL} attach a Sender link to the broker for each
        configured outbound topic mapping and forward DDS samples
        as transfers.
  \item \textbf{C.2.3 Receiver-Link Broker-to-Bridge}: the bridge
        \rfckw{SHALL} attach a Receiver link for each configured
        inbound topic mapping and forward received transfers as
        DDS samples.
  \item \textbf{C.2.4 Reconnect on Loss}: after an abrupt
        connection loss, the bridge \rfckw{SHALL} attempt to
        reconnect to the broker (per \S\,\ref{sec:reconnect}),
        \rfckw{SHALL} re-attach all configured links upon
        re-establishment, and \rfckw{SHALL NOT} drop unsettled
        DDS-side samples while the reconnect attempt is in
        progress. The exact pacing of reconnect attempts is
        configuration-dependent (\S\,\ref{sec:reconnect}); the
        test \rfckw{SHALL} verify successful re-attachment within
        the bridge's configured upper-bound back-off (default
        $\le$\,60\,s).
  \item \textbf{C.2.5 Settlement Pass-Through}: settlement
        decisions made by the broker \rfckw{SHALL} be propagated
        back to the originating DDS DataWriter.
  \item \textbf{C.2.6 Dual-Identity Separation}: a Bridge
        configured with broker-side SASL credentials
        \texttt{Alice} and an independently configured DDS-Security
        IdentityToken \texttt{CN=Bridge-1} \rfckw{SHALL} present
        \texttt{CN=Bridge-1} (not \texttt{Alice}) to the DDS-Security
        AccessControl plugin when forwarding samples
        (\S\,\ref{sec:bridge-dual-identity}). When DDS-Security
        is not active, the bridge \rfckw{SHALL} still keep the
        broker-side credential out of any DDS-side identity
        announcement.
  \item \textbf{C.2.7 Loop-Prevention (Self-Tag and Hop-Cap)}:
        two sub-tests, both \rfckw{SHALL} pass.
        (a)~The test publishes a DDS sample (DDS-side) whose
        outbound AMQP transfer carries a
        \texttt{dds:bridge-id} application-property already
        containing the bridge's own configured
        \texttt{bridge\_id}. The bridge \rfckw{SHALL NOT} forward
        the sample to the broker and \rfckw{SHALL} increment
        \texttt{transfers.dropped.loop} by one (observable via
        the bridge's \texttt{\$metrics} surface).
        (b)~A broker-originated transfer with
        \texttt{dds:bridge-hop} equal to
        \texttt{bridge\_hop\_cap} arrives on the bridge's inbound
        Receiver link. The bridge \rfckw{SHALL NOT} publish to
        DDS and \rfckw{SHALL} increment
        \texttt{transfers.dropped.hop-cap} by one.
  \item \textbf{C.2.8 Dual Management Surface}: the bridge's
        DDS-side \rfckw{SHALL} expose \texttt{\$catalog},
        \texttt{\$metrics}, and \texttt{\$audit} as Receiver
        targets (verified analogous to C.1.8/C.1.13/C.1.14). An
        attempt to attach a Receiver to any of those addresses
        on the \emph{upstream-broker} side of the bridge
        \rfckw{SHALL NOT} be made by the bridge itself, and any
        broker response that the bridge does not understand at
        those addresses \rfckw{SHALL} be propagated through to
        the operator without conflation with the DDS-side
        management surface.
\end{enumerate}

\section*{C.3 Codec-Profile Tests}

\begin{enumerate}
  \item \textbf{C.3.1 Primitive Round-Trip}: every primitive type
        of~\S\,\ref{sec:pim-type-mapping} (Primitive Mapping
        table) \rfckw{SHALL} round-trip without value change for
        a representative set of values including the type's
        minimum, maximum, and zero (where defined).
  \item \textbf{C.3.2 Compound Round-Trip}: nested
        list/map/array values up to depth 16 \rfckw{SHALL}
        round-trip without value change. (The depth bound matches
        the \rfckw{SHALL}-be-at-least-16 lower bound of
        \S\,\ref{sec:pim-impl-limits}.)
  \item \textbf{C.3.3 Compound DoS Cap}: an attempt to decode a
        compound value with depth exceeding the
        implementation-defined limit (per
        \S\,\ref{sec:pim-impl-limits}, between 16 and 64)
        \rfckw{SHALL} cause a decode error rather than stack
        overflow.
  \item \textbf{C.3.4 Performatives Round-Trip}: every performative
        of~\S\,\ref{sec:profile-codec}, item~3, \rfckw{SHALL}
        round-trip without field change for a representative
        instance.
  \item \textbf{C.3.5 Sections Sequencing}: a message containing
        all message-section variants in their canonical order
        (header, delivery-annotations, message-annotations,
        properties, application-properties, body, footer)
        \rfckw{SHALL} decode successfully; out-of-order section
        sequences \rfckw{SHALL} cause a decode error.
  \item \textbf{C.3.6 Vendor Symbolic Descriptors}: a described
        composite carrying the symbolic descriptor
        \texttt{"zerodds:time:1"}
        (\S\,\ref{sec:vendor-descriptor-encoding}) with body
        \texttt{(long sec, uint nanosec)} \rfckw{SHALL} encode
        and decode to a byte-identical wire form. An unknown
        \texttt{"zerodds:*"} descriptor \rfckw{SHALL} produce
        \texttt{amqp:not-implemented} per
        \S\,\ref{ch:error-conditions}.
\end{enumerate}

\section*{C.4 Codec-Lite-Profile Tests}

A Codec-Lite-Profile implementation \rfckw{SHALL} satisfy the
following subset of \S\,C.3:

\begin{enumerate}
  \item \textbf{C.4.1}: \S\,C.3.1 Primitive Round-Trip
        (mandatory; full primitive set).
  \item \textbf{C.4.2}: \S\,C.3.4 Performatives Round-Trip
        (mandatory; all nine performatives).
  \item \textbf{C.4.3}: \S\,C.3.5 Sections Sequencing
        \emph{restricted to} \texttt{header}, \texttt{properties},
        \texttt{data} (no compound-bearing variants), and
        \texttt{footer}.
  \item \textbf{C.4.4}: an attempt to decode an AMQP body
        whose top-level type-code is a compound (\textsf{0xC0},
        \textsf{0xC1}, \textsf{0xD0}, \textsf{0xD1},
        \textsf{0xE0}, \textsf{0xF0}) \rfckw{SHALL} produce
        \texttt{amqp:not-implemented} per
        \S\,\ref{ch:error-conditions}; the implementation is
        not required to decode compounds.
\end{enumerate}

\section*{C.5 Test Harness}

A conformance test harness is implementation-defined; the
location and toolchain of any reference harness are not part of
this specification. A harness suitable for the tests listed in
this annex \rfckw{SHALL} include at least one AMQP-1.0 client
implementation able to exercise C.1 and C.2 (for example, a
Python \texttt{pika} client, a C \texttt{qpid-proton} client, or
a .NET broker-vendor client such as
\texttt{Azure.Messaging.ServiceBus}). The codec tests of C.3 and
C.4 are exercised in-process by the implementation's own
unit-test suite. Implementers \rfckw{SHOULD} publish their
harness source so that conformance claims are independently
reproducible.

% ============================================================================
%  Annex D — DDS-RPC mapping (informative).
% ============================================================================
\annex{D}{normative when activated}{DDS-RPC Correlation Mapping}\label{annex:rpc-mapping}

This annex is conditionally normative: its requirements bind
implementations that activate the RPC-aware mapping by setting
\texttt{rpc\_aware = true} in the \texttt{TopicMapping} struct
of Annex~A. Implementations that leave \texttt{rpc\_aware =
false} (default) transport DDS-RPC samples opaquely in
Pass-Through mode (\S\,\ref{sec:psm-passthrough}) and are not
bound by this annex.

When \texttt{rpc\_aware = true}, the endpoint \rfckw{SHALL} lift
the DDS-RPC correlation header into AMQP message properties as
defined below.

\section*{D.1 Mapping Table}

When \texttt{rpc\_aware = true}, this annex \emph{overrides} the
default \texttt{message-id} mapping defined in
\S\,\ref{sec:psm-properties}: \texttt{message-id} carries the
DDS-RPC \texttt{requestId} instead of the per-sample
writer-GUID/sequence-number identifier. The per-sample identifier
that would otherwise occupy \texttt{message-id} is not carried
under RPC-aware mode; the AMQP \texttt{message-id} value
\rfckw{SHALL} be unique per RPC call, which the DDS-RPC
\texttt{requestId} guarantees by construction.

\begin{longtable}{p{4.5cm}p{8.0cm}}
  \toprule
  \textbf{DDS-RPC field} & \textbf{AMQP property} \\
  \midrule
  \endhead
  \texttt{relatedRequestId} & \texttt{correlation-id} (uuid or string) \\
  \texttt{requestId}        & \texttt{message-id} (uuid or string) — overrides the default mapping of \S\,\ref{sec:psm-properties} \\
  reply-Topic name           & \texttt{reply-to} (address string per \S\,\ref{sec:pim-address-resolution}) \\
  service-instance-name      & Application-Property \texttt{dds:rpc-instance} \\
  RPC-operation-name         & Application-Property \texttt{dds:rpc-operation} \\
  \bottomrule
\end{longtable}

The DDS \texttt{InstanceHandle\_t} continues to be carried on the
\texttt{dds:instance-handle} application-property
(\S\,\ref{sec:psm-app-properties}) regardless of \texttt{rpc\_aware}.

\section*{D.2 Activation}

The mapping is activated per topic by setting the
\texttt{rpc\_aware} flag in the \texttt{TopicMapping} struct
(Annex~A) to \texttt{true}. The default value is
\texttt{false}; when the flag is unset, DDS-RPC samples
\rfckw{SHALL} be transported in Pass-Through mode and the
correlation header remains opaque inside the XCDR2 body.

\section*{D.3 Scope of this Annex}

This annex does not define a wire-level DDS-RPC bridge; it only
suggests the property names that an RPC-aware deployment
\rfckw{SHOULD} use to permit AMQP-only consumers to participate
in DDS-RPC correlation. A full DDS-RPC-over-AMQP profile is the
subject of a future, separate specification.

\section*{D.4 Reply Validation}\label{sec:rpc-reply-validation}

When \texttt{rpc\_aware = true}, the implementation accepts AMQP
replies on behalf of an outstanding DDS-RPC caller. AMQP clients
that originate replies are not required to be DDS-aware and
\rfckw{MAY} produce replies that lack the AMQP correlation
metadata expected by Annex~D. The rules in this section bound
the failure surface so that a missing or malformed reply never
blocks the originating DDS-RPC caller indefinitely and never
mis-routes a reply to an unrelated caller.

\subsection*{D.4.1 Reply Acceptance Criteria}

A reply received on the AMQP \texttt{reply-to} address is
\rfckw{REQUIRED} to satisfy all of the following before the
implementation \rfckw{MAY} surface it to the DDS-RPC caller:

\begin{enumerate}
  \item The reply \rfckw{SHALL} carry an AMQP
        \texttt{correlation-id}.
  \item The \texttt{correlation-id} \rfckw{SHALL} match an
        outstanding RPC call that the implementation issued and
        for which it has not yet observed a reply or a timeout.
  \item The reply payload \rfckw{SHALL} decode according to the
        body-encoding mode declared on the reply topic's
        \texttt{TopicMapping}: \texttt{MODE\_PASSTHROUGH} requires
        an XCDR2 encoding of the IDL reply type per
        \omgref{OMG DDS-RPC}{1.0}; \texttt{MODE\_JSON} requires a
        JSON document conformant to \S\,\ref{sec:psm-json}, which
        the bridge \rfckw{SHALL} then decode into the IDL reply
        type; \texttt{MODE\_AMQP\_NATIVE} requires a described
        composite per \S\,\ref{sec:pim-composite-mapping}. Native
        AMQP-RPC clients that cannot produce IDL-conformant
        replies \rfckw{SHALL} be served by a reply topic
        configured for the body-encoding mode the client emits.
\end{enumerate}

\subsection*{D.4.2 Failure Dispositions}

\begin{longtable}{p{4.6cm}p{7.4cm}}
  \toprule
  \textbf{Reply condition} & \textbf{Bridge action} \\
  \midrule
  \endhead
  \texttt{correlation-id} absent &
        Reply \rfckw{SHALL} be rejected with AMQP error
        \texttt{amqp:precondition-failed}; sample
        \rfckw{SHALL} be discarded; metric
        \texttt{transfers.dropped.malformed-reply}
        \rfckw{SHALL} be incremented.
        The originating caller continues to wait until its
        per-call timeout expires (D.4.3). \\
  \texttt{correlation-id} does not match any outstanding call &
        Reply \rfckw{SHALL} be discarded silently; metric
        \texttt{transfers.dropped.malformed-reply}
        \rfckw{SHALL} be incremented. The reply \rfckw{SHALL
        NOT} be surfaced to any DDS-RPC caller. \\
  Reply payload fails IDL decode &
        Disposition \texttt{amqp:decode-error}; the originating
        DDS-RPC caller \rfckw{SHALL} receive a DDS return code of
        \texttt{RETCODE\_BAD\_PARAMETER}; metric
        \texttt{errors.decode} \rfckw{SHALL} be incremented. \\
  Reply received after \texttt{rpc\_timeout\_ms} elapsed &
        Reply \rfckw{SHALL} be discarded silently; the call has
        already terminated with \texttt{RETCODE\_TIMEOUT}; metric
        \texttt{transfers.dropped.malformed-reply} \rfckw{SHALL}
        be incremented. \\
  \bottomrule
\end{longtable}

\subsection*{D.4.3 Per-Call Timeout}

For every outstanding RPC call the implementation \rfckw{SHALL}
maintain a deadline equal to the call-issue time plus
\texttt{rpc\_timeout\_ms} (Annex~A,
\texttt{TopicMapping.rpc\_timeout\_ms}, default 30000).
When the deadline elapses without a reply that satisfies D.4.1:

\begin{itemize}
  \item the originating DDS-RPC caller \rfckw{SHALL} receive a
        DDS return code of \texttt{RETCODE\_TIMEOUT};
  \item metric \texttt{rpc.calls.timed-out} \rfckw{SHALL} be
        incremented;
  \item the corresponding entry in the implementation's
        outstanding-calls table \rfckw{SHALL} be removed so that
        a late-arriving reply with the same
        \texttt{correlation-id} is treated as no-match per the
        second row of D.4.2.
\end{itemize}

\subsection*{D.4.4 Non-Blocking Guarantee}

The implementation \rfckw{SHALL NOT} block any DDS thread or
AMQP I/O thread on a pending RPC reply. Reply validation,
correlation lookup, and timeout enforcement
\rfckw{SHALL} be performed without acquiring any lock that could
also be acquired by the DDS publish or AMQP receive paths.
Implementations \rfckw{MAY} bound the size of the
outstanding-calls table; when the bound is reached, additional
RPC-aware call attempts \rfckw{SHALL} fail synchronously with
\texttt{RETCODE\_OUT\_OF\_RESOURCES} and \rfckw{SHALL NOT} be
queued.

\end{document}
